\section{hash-binary-stack 深度题卡}

\subsection{前缀和与哈希表}
下面的题卡都按同一条链路展开：题目模型、暴力解、瓶颈、优化依据、不变量、C++ 实现、边界和扩展。格式统一是为了训练面试表达，而每张卡的关键观察和失效条件都要单独复盘。

\subsection{1 两数之和}

\textbf{题目模型。} 当前元素只需要寻找唯一补数，历史值到下标的映射就是最小状态。这题归入“前缀和与哈希表”主线，因为它的核心不是某个语法技巧，而是如何把输入状态组织成可维护的结构。读题时先把对象、关系、约束和输出目标分开：对象决定容器，关系决定边或区间，约束决定能否单调推进，输出目标决定保存最值、计数、路径还是判定结果。先查再插入，避免同一个下标被用两次；若数组有序才考虑双指针。

\textbf{暴力解和瓶颈。} 暴力做法枚举所有区间，直接求区间和或重新比较频次。即使用前缀和把求和降到常数，仍可能保留枚举左右端点的平方复杂度。 对于本题，先写暴力解的价值在于看清状态空间：哪些候选会被枚举，哪些信息会被重复计算，哪些检查其实只依赖少量历史状态。瓶颈在于当前右端需要寻找很多历史左端。若这些历史状态能被同一个键表示，就不该逐个扫描，而应把历史状态压进哈希表。 如果面试中直接跳到优化而说不清暴力慢在哪里，后续复杂度分析就会显得像背模板。

\textbf{优化依据。} 前缀和把区间问题改写成两个前缀状态的差；哈希表把寻找匹配前缀改成查表。频次表、最早位置表、最近位置表对应不同目标：计数、最长区间、最短区间不能混用。 本题的优化要抓住“先查再插入，避免同一个下标被用两次；若数组有序才考虑双指针。”这一点，把原来重复扫描或重复搜索的部分改成增量维护。优化以后，每次循环都应该只处理新增状态、删除过期状态或合并两个已经稳定的子答案，而不是重新审查全部输入。若题目条件发生变化，必须重新检查这个依据是否还成立。

\textbf{正确性不变量。} 哈希表的不变量通常是当前下标之前的历史前缀信息。先查询还是先更新不是风格问题，而是决定是否允许空区间和是否重复使用当前前缀。 用这道题复盘时，可以把不变量写成三句话：已经处理的部分满足什么，尚未处理的部分还保留什么可能性，当前操作为什么不会丢掉最优答案或合法答案。正确性来自代数等价。若两个前缀满足差值、同余或计数差要求，则它们之间的区间正好满足题意；反过来任意合法区间也对应这样一对前缀。 证明不需要很形式化，但必须能排除反例；如果你说“感觉这样选最好”，就还没有完成算法说明。

\textbf{C++ 实现要点。} 实现时用 \code{find} 判断是否存在，用 \code{operator[]} 只在确认要插入或累加时使用。总和可能溢出时使用 \code{long long}；模运算要处理负数余数归一化。本题还要特别注意：重复数字、目标为偶数且补数等于当前值、没有答案时返回空结果。写代码时先把边界条件放在函数开头或循环不变量里，而不是用很多补丁式 if 修修补补。容器选择要和访问模式一致：需要顺序就别用无序表，需要极值就别完整排序，需要历史计数就别只保存布尔值。

\textbf{复杂度和边界。} 时间复杂度要按真实输入维度说明，不能把字符串长度、图边数、值域大小都混成一个 n。空间复杂度也要区分辅助结构和输出本身。测试时至少覆盖：最小输入、没有答案、答案在边界、重复元素或重复状态、以及会让暴力解最慢的规模输入。对于这题，边界复盘重点是：重复数字、目标为偶数且补数等于当前值、没有答案时返回空结果。

\textbf{扩展和实验。} 若改成返回所有不重复数对，要先排序或用频次表处理重复。实验要覆盖从下标零开始的答案、目标为零、负数、重复前缀、余数相同但长度不足、哈希表保留第一次出现还是累加次数的差异。实验报告建议记录三件事：暴力版本在什么规模开始变慢，优化版本减少了哪类重复工作，哪一个边界用例最容易打破错误实现。这样一张题卡才算从“看过题解”变成“能够迁移”。

\subsection{49 字母异位词分组}

\textbf{题目模型。} 异位词共享同一个规范表示，可以是排序字符串或字符计数。这题归入“前缀和与哈希表”主线，因为它的核心不是某个语法技巧，而是如何把输入状态组织成可维护的结构。读题时先把对象、关系、约束和输出目标分开：对象决定容器，关系决定边或区间，约束决定能否单调推进，输出目标决定保存最值、计数、路径还是判定结果。哈希表按规范键分桶，避免两两比较字符串。

\textbf{暴力解和瓶颈。} 暴力做法枚举所有区间，直接求区间和或重新比较频次。即使用前缀和把求和降到常数，仍可能保留枚举左右端点的平方复杂度。 对于本题，先写暴力解的价值在于看清状态空间：哪些候选会被枚举，哪些信息会被重复计算，哪些检查其实只依赖少量历史状态。瓶颈在于当前右端需要寻找很多历史左端。若这些历史状态能被同一个键表示，就不该逐个扫描，而应把历史状态压进哈希表。 如果面试中直接跳到优化而说不清暴力慢在哪里，后续复杂度分析就会显得像背模板。

\textbf{优化依据。} 前缀和把区间问题改写成两个前缀状态的差；哈希表把寻找匹配前缀改成查表。频次表、最早位置表、最近位置表对应不同目标：计数、最长区间、最短区间不能混用。 本题的优化要抓住“哈希表按规范键分桶，避免两两比较字符串。”这一点，把原来重复扫描或重复搜索的部分改成增量维护。优化以后，每次循环都应该只处理新增状态、删除过期状态或合并两个已经稳定的子答案，而不是重新审查全部输入。若题目条件发生变化，必须重新检查这个依据是否还成立。

\textbf{正确性不变量。} 哈希表的不变量通常是当前下标之前的历史前缀信息。先查询还是先更新不是风格问题，而是决定是否允许空区间和是否重复使用当前前缀。 用这道题复盘时，可以把不变量写成三句话：已经处理的部分满足什么，尚未处理的部分还保留什么可能性，当前操作为什么不会丢掉最优答案或合法答案。正确性来自代数等价。若两个前缀满足差值、同余或计数差要求，则它们之间的区间正好满足题意；反过来任意合法区间也对应这样一对前缀。 证明不需要很形式化，但必须能排除反例；如果你说“感觉这样选最好”，就还没有完成算法说明。

\textbf{C++ 实现要点。} 实现时用 \code{find} 判断是否存在，用 \code{operator[]} 只在确认要插入或累加时使用。总和可能溢出时使用 \code{long long}；模运算要处理负数余数归一化。本题还要特别注意：空字符串、重复单词、字符集不止小写字母、计数键必须无歧义。写代码时先把边界条件放在函数开头或循环不变量里，而不是用很多补丁式 if 修修补补。容器选择要和访问模式一致：需要顺序就别用无序表，需要极值就别完整排序，需要历史计数就别只保存布尔值。

\textbf{复杂度和边界。} 时间复杂度要按真实输入维度说明，不能把字符串长度、图边数、值域大小都混成一个 n。空间复杂度也要区分辅助结构和输出本身。测试时至少覆盖：最小输入、没有答案、答案在边界、重复元素或重复状态、以及会让暴力解最慢的规模输入。对于这题，边界复盘重点是：空字符串、重复单词、字符集不止小写字母、计数键必须无歧义。

\textbf{扩展和实验。} 若字符串很长且字符集固定，计数键通常比排序键更稳定。实验要覆盖从下标零开始的答案、目标为零、负数、重复前缀、余数相同但长度不足、哈希表保留第一次出现还是累加次数的差异。实验报告建议记录三件事：暴力版本在什么规模开始变慢，优化版本减少了哪类重复工作，哪一个边界用例最容易打破错误实现。这样一张题卡才算从“看过题解”变成“能够迁移”。

\subsection{128 最长连续序列}

\textbf{题目模型。} 连续段只需要从没有前驱的数开始扩展。这题归入“前缀和与哈希表”主线，因为它的核心不是某个语法技巧，而是如何把输入状态组织成可维护的结构。读题时先把对象、关系、约束和输出目标分开：对象决定容器，关系决定边或区间，约束决定能否单调推进，输出目标决定保存最值、计数、路径还是判定结果。哈希集合判断 x 减一 是否存在，跳过非起点避免重复扫描。

\textbf{暴力解和瓶颈。} 暴力做法枚举所有区间，直接求区间和或重新比较频次。即使用前缀和把求和降到常数，仍可能保留枚举左右端点的平方复杂度。 对于本题，先写暴力解的价值在于看清状态空间：哪些候选会被枚举，哪些信息会被重复计算，哪些检查其实只依赖少量历史状态。瓶颈在于当前右端需要寻找很多历史左端。若这些历史状态能被同一个键表示，就不该逐个扫描，而应把历史状态压进哈希表。 如果面试中直接跳到优化而说不清暴力慢在哪里，后续复杂度分析就会显得像背模板。

\textbf{优化依据。} 前缀和把区间问题改写成两个前缀状态的差；哈希表把寻找匹配前缀改成查表。频次表、最早位置表、最近位置表对应不同目标：计数、最长区间、最短区间不能混用。 本题的优化要抓住“哈希集合判断 x 减一 是否存在，跳过非起点避免重复扫描。”这一点，把原来重复扫描或重复搜索的部分改成增量维护。优化以后，每次循环都应该只处理新增状态、删除过期状态或合并两个已经稳定的子答案，而不是重新审查全部输入。若题目条件发生变化，必须重新检查这个依据是否还成立。

\textbf{正确性不变量。} 哈希表的不变量通常是当前下标之前的历史前缀信息。先查询还是先更新不是风格问题，而是决定是否允许空区间和是否重复使用当前前缀。 用这道题复盘时，可以把不变量写成三句话：已经处理的部分满足什么，尚未处理的部分还保留什么可能性，当前操作为什么不会丢掉最优答案或合法答案。正确性来自代数等价。若两个前缀满足差值、同余或计数差要求，则它们之间的区间正好满足题意；反过来任意合法区间也对应这样一对前缀。 证明不需要很形式化，但必须能排除反例；如果你说“感觉这样选最好”，就还没有完成算法说明。

\textbf{C++ 实现要点。} 实现时用 \code{find} 判断是否存在，用 \code{operator[]} 只在确认要插入或累加时使用。总和可能溢出时使用 \code{long long}；模运算要处理负数余数归一化。本题还要特别注意：重复元素、负数、空数组、整数边界。写代码时先把边界条件放在函数开头或循环不变量里，而不是用很多补丁式 if 修修补补。容器选择要和访问模式一致：需要顺序就别用无序表，需要极值就别完整排序，需要历史计数就别只保存布尔值。

\textbf{复杂度和边界。} 时间复杂度要按真实输入维度说明，不能把字符串长度、图边数、值域大小都混成一个 n。空间复杂度也要区分辅助结构和输出本身。测试时至少覆盖：最小输入、没有答案、答案在边界、重复元素或重复状态、以及会让暴力解最慢的规模输入。对于这题，边界复盘重点是：重复元素、负数、空数组、整数边界。

\textbf{扩展和实验。} 若数据流动态插入，可以用并查集或区间合并维护长度。实验要覆盖从下标零开始的答案、目标为零、负数、重复前缀、余数相同但长度不足、哈希表保留第一次出现还是累加次数的差异。实验报告建议记录三件事：暴力版本在什么规模开始变慢，优化版本减少了哪类重复工作，哪一个边界用例最容易打破错误实现。这样一张题卡才算从“看过题解”变成“能够迁移”。

\subsection{202 快乐数}

\textbf{题目模型。} 反复变换数字会进入一或循环，状态空间最终有限。这题归入“前缀和与哈希表”主线，因为它的核心不是某个语法技巧，而是如何把输入状态组织成可维护的结构。读题时先把对象、关系、约束和输出目标分开：对象决定容器，关系决定边或区间，约束决定能否单调推进，输出目标决定保存最值、计数、路径还是判定结果。用集合检测重复状态，或用快慢指针看成函数图。

\textbf{暴力解和瓶颈。} 暴力做法枚举所有区间，直接求区间和或重新比较频次。即使用前缀和把求和降到常数，仍可能保留枚举左右端点的平方复杂度。 对于本题，先写暴力解的价值在于看清状态空间：哪些候选会被枚举，哪些信息会被重复计算，哪些检查其实只依赖少量历史状态。瓶颈在于当前右端需要寻找很多历史左端。若这些历史状态能被同一个键表示，就不该逐个扫描，而应把历史状态压进哈希表。 如果面试中直接跳到优化而说不清暴力慢在哪里，后续复杂度分析就会显得像背模板。

\textbf{优化依据。} 前缀和把区间问题改写成两个前缀状态的差；哈希表把寻找匹配前缀改成查表。频次表、最早位置表、最近位置表对应不同目标：计数、最长区间、最短区间不能混用。 本题的优化要抓住“用集合检测重复状态，或用快慢指针看成函数图。”这一点，把原来重复扫描或重复搜索的部分改成增量维护。优化以后，每次循环都应该只处理新增状态、删除过期状态或合并两个已经稳定的子答案，而不是重新审查全部输入。若题目条件发生变化，必须重新检查这个依据是否还成立。

\textbf{正确性不变量。} 哈希表的不变量通常是当前下标之前的历史前缀信息。先查询还是先更新不是风格问题，而是决定是否允许空区间和是否重复使用当前前缀。 用这道题复盘时，可以把不变量写成三句话：已经处理的部分满足什么，尚未处理的部分还保留什么可能性，当前操作为什么不会丢掉最优答案或合法答案。正确性来自代数等价。若两个前缀满足差值、同余或计数差要求，则它们之间的区间正好满足题意；反过来任意合法区间也对应这样一对前缀。 证明不需要很形式化，但必须能排除反例；如果你说“感觉这样选最好”，就还没有完成算法说明。

\textbf{C++ 实现要点。} 实现时用 \code{find} 判断是否存在，用 \code{operator[]} 只在确认要插入或累加时使用。总和可能溢出时使用 \code{long long}；模运算要处理负数余数归一化。本题还要特别注意：输入为一、进入循环、位平方和函数、负数不在题意内。写代码时先把边界条件放在函数开头或循环不变量里，而不是用很多补丁式 if 修修补补。容器选择要和访问模式一致：需要顺序就别用无序表，需要极值就别完整排序，需要历史计数就别只保存布尔值。

\textbf{复杂度和边界。} 时间复杂度要按真实输入维度说明，不能把字符串长度、图边数、值域大小都混成一个 n。空间复杂度也要区分辅助结构和输出本身。测试时至少覆盖：最小输入、没有答案、答案在边界、重复元素或重复状态、以及会让暴力解最慢的规模输入。对于这题，边界复盘重点是：输入为一、进入循环、位平方和函数、负数不在题意内。

\textbf{扩展和实验。} 快慢指针版本能训练把数值迭代看成链表。实验要覆盖从下标零开始的答案、目标为零、负数、重复前缀、余数相同但长度不足、哈希表保留第一次出现还是累加次数的差异。实验报告建议记录三件事：暴力版本在什么规模开始变慢，优化版本减少了哪类重复工作，哪一个边界用例最容易打破错误实现。这样一张题卡才算从“看过题解”变成“能够迁移”。

\subsection{217 存在重复元素}

\textbf{题目模型。} 判断是否出现过是哈希集合最直接的职责。这题归入“前缀和与哈希表”主线，因为它的核心不是某个语法技巧，而是如何把输入状态组织成可维护的结构。读题时先把对象、关系、约束和输出目标分开：对象决定容器，关系决定边或区间，约束决定能否单调推进，输出目标决定保存最值、计数、路径还是判定结果。扫描时若插入前已经存在则立即返回真。

\textbf{暴力解和瓶颈。} 暴力做法枚举所有区间，直接求区间和或重新比较频次。即使用前缀和把求和降到常数，仍可能保留枚举左右端点的平方复杂度。 对于本题，先写暴力解的价值在于看清状态空间：哪些候选会被枚举，哪些信息会被重复计算，哪些检查其实只依赖少量历史状态。瓶颈在于当前右端需要寻找很多历史左端。若这些历史状态能被同一个键表示，就不该逐个扫描，而应把历史状态压进哈希表。 如果面试中直接跳到优化而说不清暴力慢在哪里，后续复杂度分析就会显得像背模板。

\textbf{优化依据。} 前缀和把区间问题改写成两个前缀状态的差；哈希表把寻找匹配前缀改成查表。频次表、最早位置表、最近位置表对应不同目标：计数、最长区间、最短区间不能混用。 本题的优化要抓住“扫描时若插入前已经存在则立即返回真。”这一点，把原来重复扫描或重复搜索的部分改成增量维护。优化以后，每次循环都应该只处理新增状态、删除过期状态或合并两个已经稳定的子答案，而不是重新审查全部输入。若题目条件发生变化，必须重新检查这个依据是否还成立。

\textbf{正确性不变量。} 哈希表的不变量通常是当前下标之前的历史前缀信息。先查询还是先更新不是风格问题，而是决定是否允许空区间和是否重复使用当前前缀。 用这道题复盘时，可以把不变量写成三句话：已经处理的部分满足什么，尚未处理的部分还保留什么可能性，当前操作为什么不会丢掉最优答案或合法答案。正确性来自代数等价。若两个前缀满足差值、同余或计数差要求，则它们之间的区间正好满足题意；反过来任意合法区间也对应这样一对前缀。 证明不需要很形式化，但必须能排除反例；如果你说“感觉这样选最好”，就还没有完成算法说明。

\textbf{C++ 实现要点。} 实现时用 \code{find} 判断是否存在，用 \code{operator[]} 只在确认要插入或累加时使用。总和可能溢出时使用 \code{long long}；模运算要处理负数余数归一化。本题还要特别注意：空数组、全唯一、重复在末尾、值域很小。写代码时先把边界条件放在函数开头或循环不变量里，而不是用很多补丁式 if 修修补补。容器选择要和访问模式一致：需要顺序就别用无序表，需要极值就别完整排序，需要历史计数就别只保存布尔值。

\textbf{复杂度和边界。} 时间复杂度要按真实输入维度说明，不能把字符串长度、图边数、值域大小都混成一个 n。空间复杂度也要区分辅助结构和输出本身。测试时至少覆盖：最小输入、没有答案、答案在边界、重复元素或重复状态、以及会让暴力解最慢的规模输入。对于这题，边界复盘重点是：空数组、全唯一、重复在末尾、值域很小。

\textbf{扩展和实验。} 若允许修改输入，排序后看相邻可以降低额外空间。实验要覆盖从下标零开始的答案、目标为零、负数、重复前缀、余数相同但长度不足、哈希表保留第一次出现还是累加次数的差异。实验报告建议记录三件事：暴力版本在什么规模开始变慢，优化版本减少了哪类重复工作，哪一个边界用例最容易打破错误实现。这样一张题卡才算从“看过题解”变成“能够迁移”。

\subsection{560 和为 K 的子数组}

\textbf{题目模型。} 当前前缀和减去目标值，就是需要匹配的历史前缀。这题归入“前缀和与哈希表”主线，因为它的核心不是某个语法技巧，而是如何把输入状态组织成可维护的结构。读题时先把对象、关系、约束和输出目标分开：对象决定容器，关系决定边或区间，约束决定能否单调推进，输出目标决定保存最值、计数、路径还是判定结果。哈希表保存前缀和出现次数，先查询再更新。

\textbf{暴力解和瓶颈。} 暴力做法枚举所有区间，直接求区间和或重新比较频次。即使用前缀和把求和降到常数，仍可能保留枚举左右端点的平方复杂度。 对于本题，先写暴力解的价值在于看清状态空间：哪些候选会被枚举，哪些信息会被重复计算，哪些检查其实只依赖少量历史状态。瓶颈在于当前右端需要寻找很多历史左端。若这些历史状态能被同一个键表示，就不该逐个扫描，而应把历史状态压进哈希表。 如果面试中直接跳到优化而说不清暴力慢在哪里，后续复杂度分析就会显得像背模板。

\textbf{优化依据。} 前缀和把区间问题改写成两个前缀状态的差；哈希表把寻找匹配前缀改成查表。频次表、最早位置表、最近位置表对应不同目标：计数、最长区间、最短区间不能混用。 本题的优化要抓住“哈希表保存前缀和出现次数，先查询再更新。”这一点，把原来重复扫描或重复搜索的部分改成增量维护。优化以后，每次循环都应该只处理新增状态、删除过期状态或合并两个已经稳定的子答案，而不是重新审查全部输入。若题目条件发生变化，必须重新检查这个依据是否还成立。

\textbf{正确性不变量。} 哈希表的不变量通常是当前下标之前的历史前缀信息。先查询还是先更新不是风格问题，而是决定是否允许空区间和是否重复使用当前前缀。 用这道题复盘时，可以把不变量写成三句话：已经处理的部分满足什么，尚未处理的部分还保留什么可能性，当前操作为什么不会丢掉最优答案或合法答案。正确性来自代数等价。若两个前缀满足差值、同余或计数差要求，则它们之间的区间正好满足题意；反过来任意合法区间也对应这样一对前缀。 证明不需要很形式化，但必须能排除反例；如果你说“感觉这样选最好”，就还没有完成算法说明。

\textbf{C++ 实现要点。} 实现时用 \code{find} 判断是否存在，用 \code{operator[]} 只在确认要插入或累加时使用。总和可能溢出时使用 \code{long long}；模运算要处理负数余数归一化。本题还要特别注意：负数、目标为零、从下标零开始的区间、重复前缀。写代码时先把边界条件放在函数开头或循环不变量里，而不是用很多补丁式 if 修修补补。容器选择要和访问模式一致：需要顺序就别用无序表，需要极值就别完整排序，需要历史计数就别只保存布尔值。

\textbf{复杂度和边界。} 时间复杂度要按真实输入维度说明，不能把字符串长度、图边数、值域大小都混成一个 n。空间复杂度也要区分辅助结构和输出本身。测试时至少覆盖：最小输入、没有答案、答案在边界、重复元素或重复状态、以及会让暴力解最慢的规模输入。对于这题，边界复盘重点是：负数、目标为零、从下标零开始的区间、重复前缀。

\textbf{扩展和实验。} 若要求最长区间，哈希表应保存最早位置而不是次数。实验要覆盖从下标零开始的答案、目标为零、负数、重复前缀、余数相同但长度不足、哈希表保留第一次出现还是累加次数的差异。实验报告建议记录三件事：暴力版本在什么规模开始变慢，优化版本减少了哪类重复工作，哪一个边界用例最容易打破错误实现。这样一张题卡才算从“看过题解”变成“能够迁移”。

\subsection{二分查找与答案空间}
下面的题卡都按同一条链路展开：题目模型、暴力解、瓶颈、优化依据、不变量、C++ 实现、边界和扩展。格式统一是为了训练面试表达，而每张卡的关键观察和失效条件都要单独复盘。

\subsection{33 搜索旋转排序数组}

\textbf{题目模型。} 旋转数组每次二分至少有一半保持有序。这题归入“二分查找与答案空间”主线，因为它的核心不是某个语法技巧，而是如何把输入状态组织成可维护的结构。读题时先把对象、关系、约束和输出目标分开：对象决定容器，关系决定边或区间，约束决定能否单调推进，输出目标决定保存最值、计数、路径还是判定结果。判断哪一半有序，再看目标是否落在有序半边范围内。

\textbf{暴力解和瓶颈。} 暴力做法通常线性扫描数组或线性枚举答案。它的问题不是不会找到答案，而是没有利用有序性或可行性的单调边界。 对于本题，先写暴力解的价值在于看清状态空间：哪些候选会被枚举，哪些信息会被重复计算，哪些检查其实只依赖少量历史状态。慢点在于每次只排除一个候选。二分能成立，是因为一次比较可以排除一半下标，或者一次可行性检查可以排除一半答案范围。 如果面试中直接跳到优化而说不清暴力慢在哪里，后续复杂度分析就会显得像背模板。

\textbf{优化依据。} 二分前必须先定义谓词：某个位置是否已经满足边界条件，或者某个答案是否可行。数组二分找的是边界，答案二分找的是最小可行值或最大可行值。 本题的优化要抓住“判断哪一半有序，再看目标是否落在有序半边范围内。”这一点，把原来重复扫描或重复搜索的部分改成增量维护。优化以后，每次循环都应该只处理新增状态、删除过期状态或合并两个已经稳定的子答案，而不是重新审查全部输入。若题目条件发生变化，必须重新检查这个依据是否还成立。

\textbf{正确性不变量。} 建议固定左闭右开区间。循环中始终维护答案在 \code{[left, right)} 内，左侧一定不可行或小于目标，右侧外部已经被排除。 用这道题复盘时，可以把不变量写成三句话：已经处理的部分满足什么，尚未处理的部分还保留什么可能性，当前操作为什么不会丢掉最优答案或合法答案。证明二分不是证明循环会停，而是证明被排除的一半确实不含答案。答案空间二分还要证明可行性谓词单调。 证明不需要很形式化，但必须能排除反例；如果你说“感觉这样选最好”，就还没有完成算法说明。

\textbf{C++ 实现要点。} 中点写成 \code{left + (right - left) / 2}。判断平方、乘法、总量时避免溢出。用 \code{std::lower_bound} 能表达清楚时优先用它，但面试要能手写同等语义。本题还要特别注意：未旋转、长度一、目标不存在、边界等号、存在重复元素时退化。写代码时先把边界条件放在函数开头或循环不变量里，而不是用很多补丁式 if 修修补补。容器选择要和访问模式一致：需要顺序就别用无序表，需要极值就别完整排序，需要历史计数就别只保存布尔值。

\textbf{复杂度和边界。} 时间复杂度要按真实输入维度说明，不能把字符串长度、图边数、值域大小都混成一个 n。空间复杂度也要区分辅助结构和输出本身。测试时至少覆盖：最小输入、没有答案、答案在边界、重复元素或重复状态、以及会让暴力解最慢的规模输入。对于这题，边界复盘重点是：未旋转、长度一、目标不存在、边界等号、存在重复元素时退化。

\textbf{扩展和实验。} 若重复元素很多，需要在无法判断有序半边时收缩边界。实验准备长度为零、一、二的数组，目标小于全部元素、大于全部元素、重复元素边界、旋转数组、不可行答案和刚好可行答案。实验报告建议记录三件事：暴力版本在什么规模开始变慢，优化版本减少了哪类重复工作，哪一个边界用例最容易打破错误实现。这样一张题卡才算从“看过题解”变成“能够迁移”。

\subsection{34 排序数组中查找首尾位置}

\textbf{题目模型。} 首尾位置可以拆成两个边界：第一个大于等于和第一个大于。这题归入“二分查找与答案空间”主线，因为它的核心不是某个语法技巧，而是如何把输入状态组织成可维护的结构。读题时先把对象、关系、约束和输出目标分开：对象决定容器，关系决定边或区间，约束决定能否单调推进，输出目标决定保存最值、计数、路径还是判定结果。不要找到一个目标后向两边线性扩展，否则重复元素很多时退化。

\textbf{暴力解和瓶颈。} 暴力做法通常线性扫描数组或线性枚举答案。它的问题不是不会找到答案，而是没有利用有序性或可行性的单调边界。 对于本题，先写暴力解的价值在于看清状态空间：哪些候选会被枚举，哪些信息会被重复计算，哪些检查其实只依赖少量历史状态。慢点在于每次只排除一个候选。二分能成立，是因为一次比较可以排除一半下标，或者一次可行性检查可以排除一半答案范围。 如果面试中直接跳到优化而说不清暴力慢在哪里，后续复杂度分析就会显得像背模板。

\textbf{优化依据。} 二分前必须先定义谓词：某个位置是否已经满足边界条件，或者某个答案是否可行。数组二分找的是边界，答案二分找的是最小可行值或最大可行值。 本题的优化要抓住“不要找到一个目标后向两边线性扩展，否则重复元素很多时退化。”这一点，把原来重复扫描或重复搜索的部分改成增量维护。优化以后，每次循环都应该只处理新增状态、删除过期状态或合并两个已经稳定的子答案，而不是重新审查全部输入。若题目条件发生变化，必须重新检查这个依据是否还成立。

\textbf{正确性不变量。} 建议固定左闭右开区间。循环中始终维护答案在 \code{[left, right)} 内，左侧一定不可行或小于目标，右侧外部已经被排除。 用这道题复盘时，可以把不变量写成三句话：已经处理的部分满足什么，尚未处理的部分还保留什么可能性，当前操作为什么不会丢掉最优答案或合法答案。证明二分不是证明循环会停，而是证明被排除的一半确实不含答案。答案空间二分还要证明可行性谓词单调。 证明不需要很形式化，但必须能排除反例；如果你说“感觉这样选最好”，就还没有完成算法说明。

\textbf{C++ 实现要点。} 中点写成 \code{left + (right - left) / 2}。判断平方、乘法、总量时避免溢出。用 \code{std::lower_bound} 能表达清楚时优先用它，但面试要能手写同等语义。本题还要特别注意：目标不存在、目标在开头或结尾、数组为空、全是目标。写代码时先把边界条件放在函数开头或循环不变量里，而不是用很多补丁式 if 修修补补。容器选择要和访问模式一致：需要顺序就别用无序表，需要极值就别完整排序，需要历史计数就别只保存布尔值。

\textbf{复杂度和边界。} 时间复杂度要按真实输入维度说明，不能把字符串长度、图边数、值域大小都混成一个 n。空间复杂度也要区分辅助结构和输出本身。测试时至少覆盖：最小输入、没有答案、答案在边界、重复元素或重复状态、以及会让暴力解最慢的规模输入。对于这题，边界复盘重点是：目标不存在、目标在开头或结尾、数组为空、全是目标。

\textbf{扩展和实验。} 这个模型能迁移到计数某值出现次数和插入区间边界。实验准备长度为零、一、二的数组，目标小于全部元素、大于全部元素、重复元素边界、旋转数组、不可行答案和刚好可行答案。实验报告建议记录三件事：暴力版本在什么规模开始变慢，优化版本减少了哪类重复工作，哪一个边界用例最容易打破错误实现。这样一张题卡才算从“看过题解”变成“能够迁移”。

\subsection{35 搜索插入位置}

\textbf{题目模型。} 题目等价于寻找第一个大于等于目标的位置。这题归入“二分查找与答案空间”主线，因为它的核心不是某个语法技巧，而是如何把输入状态组织成可维护的结构。读题时先把对象、关系、约束和输出目标分开：对象决定容器，关系决定边或区间，约束决定能否单调推进，输出目标决定保存最值、计数、路径还是判定结果。统一用 lower bound 语义，不在相等时提前返回也能保持边界一致。

\textbf{暴力解和瓶颈。} 暴力做法通常线性扫描数组或线性枚举答案。它的问题不是不会找到答案，而是没有利用有序性或可行性的单调边界。 对于本题，先写暴力解的价值在于看清状态空间：哪些候选会被枚举，哪些信息会被重复计算，哪些检查其实只依赖少量历史状态。慢点在于每次只排除一个候选。二分能成立，是因为一次比较可以排除一半下标，或者一次可行性检查可以排除一半答案范围。 如果面试中直接跳到优化而说不清暴力慢在哪里，后续复杂度分析就会显得像背模板。

\textbf{优化依据。} 二分前必须先定义谓词：某个位置是否已经满足边界条件，或者某个答案是否可行。数组二分找的是边界，答案二分找的是最小可行值或最大可行值。 本题的优化要抓住“统一用 lower bound 语义，不在相等时提前返回也能保持边界一致。”这一点，把原来重复扫描或重复搜索的部分改成增量维护。优化以后，每次循环都应该只处理新增状态、删除过期状态或合并两个已经稳定的子答案，而不是重新审查全部输入。若题目条件发生变化，必须重新检查这个依据是否还成立。

\textbf{正确性不变量。} 建议固定左闭右开区间。循环中始终维护答案在 \code{[left, right)} 内，左侧一定不可行或小于目标，右侧外部已经被排除。 用这道题复盘时，可以把不变量写成三句话：已经处理的部分满足什么，尚未处理的部分还保留什么可能性，当前操作为什么不会丢掉最优答案或合法答案。证明二分不是证明循环会停，而是证明被排除的一半确实不含答案。答案空间二分还要证明可行性谓词单调。 证明不需要很形式化，但必须能排除反例；如果你说“感觉这样选最好”，就还没有完成算法说明。

\textbf{C++ 实现要点。} 中点写成 \code{left + (right - left) / 2}。判断平方、乘法、总量时避免溢出。用 \code{std::lower_bound} 能表达清楚时优先用它，但面试要能手写同等语义。本题还要特别注意：目标小于所有元素、目标大于所有元素、重复值、空数组。写代码时先把边界条件放在函数开头或循环不变量里，而不是用很多补丁式 if 修修补补。容器选择要和访问模式一致：需要顺序就别用无序表，需要极值就别完整排序，需要历史计数就别只保存布尔值。

\textbf{复杂度和边界。} 时间复杂度要按真实输入维度说明，不能把字符串长度、图边数、值域大小都混成一个 n。空间复杂度也要区分辅助结构和输出本身。测试时至少覆盖：最小输入、没有答案、答案在边界、重复元素或重复状态、以及会让暴力解最慢的规模输入。对于这题，边界复盘重点是：目标小于所有元素、目标大于所有元素、重复值、空数组。

\textbf{扩展和实验。} 手写二分时用左闭右开可直接返回 left。实验准备长度为零、一、二的数组，目标小于全部元素、大于全部元素、重复元素边界、旋转数组、不可行答案和刚好可行答案。实验报告建议记录三件事：暴力版本在什么规模开始变慢，优化版本减少了哪类重复工作，哪一个边界用例最容易打破错误实现。这样一张题卡才算从“看过题解”变成“能够迁移”。

\subsection{153 寻找旋转排序数组中的最小值}

\textbf{题目模型。} 最小值是旋转断点，右端点可帮助判断中点在哪一段。这题归入“二分查找与答案空间”主线，因为它的核心不是某个语法技巧，而是如何把输入状态组织成可维护的结构。读题时先把对象、关系、约束和输出目标分开：对象决定容器，关系决定边或区间，约束决定能否单调推进，输出目标决定保存最值、计数、路径还是判定结果。若 mid 大于 right，最小值在右侧；否则在左侧含 mid。

\textbf{暴力解和瓶颈。} 暴力做法通常线性扫描数组或线性枚举答案。它的问题不是不会找到答案，而是没有利用有序性或可行性的单调边界。 对于本题，先写暴力解的价值在于看清状态空间：哪些候选会被枚举，哪些信息会被重复计算，哪些检查其实只依赖少量历史状态。慢点在于每次只排除一个候选。二分能成立，是因为一次比较可以排除一半下标，或者一次可行性检查可以排除一半答案范围。 如果面试中直接跳到优化而说不清暴力慢在哪里，后续复杂度分析就会显得像背模板。

\textbf{优化依据。} 二分前必须先定义谓词：某个位置是否已经满足边界条件，或者某个答案是否可行。数组二分找的是边界，答案二分找的是最小可行值或最大可行值。 本题的优化要抓住“若 mid 大于 right，最小值在右侧；否则在左侧含 mid。”这一点，把原来重复扫描或重复搜索的部分改成增量维护。优化以后，每次循环都应该只处理新增状态、删除过期状态或合并两个已经稳定的子答案，而不是重新审查全部输入。若题目条件发生变化，必须重新检查这个依据是否还成立。

\textbf{正确性不变量。} 建议固定左闭右开区间。循环中始终维护答案在 \code{[left, right)} 内，左侧一定不可行或小于目标，右侧外部已经被排除。 用这道题复盘时，可以把不变量写成三句话：已经处理的部分满足什么，尚未处理的部分还保留什么可能性，当前操作为什么不会丢掉最优答案或合法答案。证明二分不是证明循环会停，而是证明被排除的一半确实不含答案。答案空间二分还要证明可行性谓词单调。 证明不需要很形式化，但必须能排除反例；如果你说“感觉这样选最好”，就还没有完成算法说明。

\textbf{C++ 实现要点。} 中点写成 \code{left + (right - left) / 2}。判断平方、乘法、总量时避免溢出。用 \code{std::lower_bound} 能表达清楚时优先用它，但面试要能手写同等语义。本题还要特别注意：未旋转、长度一、旋转一位、无重复条件。写代码时先把边界条件放在函数开头或循环不变量里，而不是用很多补丁式 if 修修补补。容器选择要和访问模式一致：需要顺序就别用无序表，需要极值就别完整排序，需要历史计数就别只保存布尔值。

\textbf{复杂度和边界。} 时间复杂度要按真实输入维度说明，不能把字符串长度、图边数、值域大小都混成一个 n。空间复杂度也要区分辅助结构和输出本身。测试时至少覆盖：最小输入、没有答案、答案在边界、重复元素或重复状态、以及会让暴力解最慢的规模输入。对于这题，边界复盘重点是：未旋转、长度一、旋转一位、无重复条件。

\textbf{扩展和实验。} 有重复时遇到相等只能收缩 right，最坏退化。实验准备长度为零、一、二的数组，目标小于全部元素、大于全部元素、重复元素边界、旋转数组、不可行答案和刚好可行答案。实验报告建议记录三件事：暴力版本在什么规模开始变慢，优化版本减少了哪类重复工作，哪一个边界用例最容易打破错误实现。这样一张题卡才算从“看过题解”变成“能够迁移”。

\subsection{154 寻找旋转排序数组中的最小值 II}

\textbf{题目模型。} 重复元素会让 mid 和 right 相等时无法判断最小值方向。这题归入“二分查找与答案空间”主线，因为它的核心不是某个语法技巧，而是如何把输入状态组织成可维护的结构。读题时先把对象、关系、约束和输出目标分开：对象决定容器，关系决定边或区间，约束决定能否单调推进，输出目标决定保存最值、计数、路径还是判定结果。相等时只能安全地丢掉一个 right，因为它不是唯一必要候选。

\textbf{暴力解和瓶颈。} 暴力做法通常线性扫描数组或线性枚举答案。它的问题不是不会找到答案，而是没有利用有序性或可行性的单调边界。 对于本题，先写暴力解的价值在于看清状态空间：哪些候选会被枚举，哪些信息会被重复计算，哪些检查其实只依赖少量历史状态。慢点在于每次只排除一个候选。二分能成立，是因为一次比较可以排除一半下标，或者一次可行性检查可以排除一半答案范围。 如果面试中直接跳到优化而说不清暴力慢在哪里，后续复杂度分析就会显得像背模板。

\textbf{优化依据。} 二分前必须先定义谓词：某个位置是否已经满足边界条件，或者某个答案是否可行。数组二分找的是边界，答案二分找的是最小可行值或最大可行值。 本题的优化要抓住“相等时只能安全地丢掉一个 right，因为它不是唯一必要候选。”这一点，把原来重复扫描或重复搜索的部分改成增量维护。优化以后，每次循环都应该只处理新增状态、删除过期状态或合并两个已经稳定的子答案，而不是重新审查全部输入。若题目条件发生变化，必须重新检查这个依据是否还成立。

\textbf{正确性不变量。} 建议固定左闭右开区间。循环中始终维护答案在 \code{[left, right)} 内，左侧一定不可行或小于目标，右侧外部已经被排除。 用这道题复盘时，可以把不变量写成三句话：已经处理的部分满足什么，尚未处理的部分还保留什么可能性，当前操作为什么不会丢掉最优答案或合法答案。证明二分不是证明循环会停，而是证明被排除的一半确实不含答案。答案空间二分还要证明可行性谓词单调。 证明不需要很形式化，但必须能排除反例；如果你说“感觉这样选最好”，就还没有完成算法说明。

\textbf{C++ 实现要点。} 中点写成 \code{left + (right - left) / 2}。判断平方、乘法、总量时避免溢出。用 \code{std::lower_bound} 能表达清楚时优先用它，但面试要能手写同等语义。本题还要特别注意：全相等、重复值跨断点、最小值出现多次。写代码时先把边界条件放在函数开头或循环不变量里，而不是用很多补丁式 if 修修补补。容器选择要和访问模式一致：需要顺序就别用无序表，需要极值就别完整排序，需要历史计数就别只保存布尔值。

\textbf{复杂度和边界。} 时间复杂度要按真实输入维度说明，不能把字符串长度、图边数、值域大小都混成一个 n。空间复杂度也要区分辅助结构和输出本身。测试时至少覆盖：最小输入、没有答案、答案在边界、重复元素或重复状态、以及会让暴力解最慢的规模输入。对于这题，边界复盘重点是：全相等、重复值跨断点、最小值出现多次。

\textbf{扩展和实验。} 这题说明二分依赖的信息不足时会退化，不是所有旋转题都严格对数。实验准备长度为零、一、二的数组，目标小于全部元素、大于全部元素、重复元素边界、旋转数组、不可行答案和刚好可行答案。实验报告建议记录三件事：暴力版本在什么规模开始变慢，优化版本减少了哪类重复工作，哪一个边界用例最容易打破错误实现。这样一张题卡才算从“看过题解”变成“能够迁移”。

\subsection{162 寻找峰值}

\textbf{题目模型。} 比较 mid 和 mid 加一 的斜率能判断某侧一定存在峰值。这题归入“二分查找与答案空间”主线，因为它的核心不是某个语法技巧，而是如何把输入状态组织成可维护的结构。读题时先把对象、关系、约束和输出目标分开：对象决定容器，关系决定边或区间，约束决定能否单调推进，输出目标决定保存最值、计数、路径还是判定结果。上升则右侧有峰，下降则左侧含 mid 有峰。

\textbf{暴力解和瓶颈。} 暴力做法通常线性扫描数组或线性枚举答案。它的问题不是不会找到答案，而是没有利用有序性或可行性的单调边界。 对于本题，先写暴力解的价值在于看清状态空间：哪些候选会被枚举，哪些信息会被重复计算，哪些检查其实只依赖少量历史状态。慢点在于每次只排除一个候选。二分能成立，是因为一次比较可以排除一半下标，或者一次可行性检查可以排除一半答案范围。 如果面试中直接跳到优化而说不清暴力慢在哪里，后续复杂度分析就会显得像背模板。

\textbf{优化依据。} 二分前必须先定义谓词：某个位置是否已经满足边界条件，或者某个答案是否可行。数组二分找的是边界，答案二分找的是最小可行值或最大可行值。 本题的优化要抓住“上升则右侧有峰，下降则左侧含 mid 有峰。”这一点，把原来重复扫描或重复搜索的部分改成增量维护。优化以后，每次循环都应该只处理新增状态、删除过期状态或合并两个已经稳定的子答案，而不是重新审查全部输入。若题目条件发生变化，必须重新检查这个依据是否还成立。

\textbf{正确性不变量。} 建议固定左闭右开区间。循环中始终维护答案在 \code{[left, right)} 内，左侧一定不可行或小于目标，右侧外部已经被排除。 用这道题复盘时，可以把不变量写成三句话：已经处理的部分满足什么，尚未处理的部分还保留什么可能性，当前操作为什么不会丢掉最优答案或合法答案。证明二分不是证明循环会停，而是证明被排除的一半确实不含答案。答案空间二分还要证明可行性谓词单调。 证明不需要很形式化，但必须能排除反例；如果你说“感觉这样选最好”，就还没有完成算法说明。

\textbf{C++ 实现要点。} 中点写成 \code{left + (right - left) / 2}。判断平方、乘法、总量时避免溢出。用 \code{std::lower_bound} 能表达清楚时优先用它，但面试要能手写同等语义。本题还要特别注意：长度一、峰在边界、多个峰、相邻不相等假设。写代码时先把边界条件放在函数开头或循环不变量里，而不是用很多补丁式 if 修修补补。容器选择要和访问模式一致：需要顺序就别用无序表，需要极值就别完整排序，需要历史计数就别只保存布尔值。

\textbf{复杂度和边界。} 时间复杂度要按真实输入维度说明，不能把字符串长度、图边数、值域大小都混成一个 n。空间复杂度也要区分辅助结构和输出本身。测试时至少覆盖：最小输入、没有答案、答案在边界、重复元素或重复状态、以及会让暴力解最慢的规模输入。对于这题，边界复盘重点是：长度一、峰在边界、多个峰、相邻不相等假设。

\textbf{扩展和实验。} 二维峰值需要按行列最大值进行更复杂二分。实验准备长度为零、一、二的数组，目标小于全部元素、大于全部元素、重复元素边界、旋转数组、不可行答案和刚好可行答案。实验报告建议记录三件事：暴力版本在什么规模开始变慢，优化版本减少了哪类重复工作，哪一个边界用例最容易打破错误实现。这样一张题卡才算从“看过题解”变成“能够迁移”。

\subsection{240 搜索二维矩阵 II}

\textbf{题目模型。} 右上角同时具有向左变小、向下变大的排除能力。这题归入“二分查找与答案空间”主线，因为它的核心不是某个语法技巧，而是如何把输入状态组织成可维护的结构。读题时先把对象、关系、约束和输出目标分开：对象决定容器，关系决定边或区间，约束决定能否单调推进，输出目标决定保存最值、计数、路径还是判定结果。当前值大于目标就左移，小于目标就下移。

\textbf{暴力解和瓶颈。} 暴力做法通常线性扫描数组或线性枚举答案。它的问题不是不会找到答案，而是没有利用有序性或可行性的单调边界。 对于本题，先写暴力解的价值在于看清状态空间：哪些候选会被枚举，哪些信息会被重复计算，哪些检查其实只依赖少量历史状态。慢点在于每次只排除一个候选。二分能成立，是因为一次比较可以排除一半下标，或者一次可行性检查可以排除一半答案范围。 如果面试中直接跳到优化而说不清暴力慢在哪里，后续复杂度分析就会显得像背模板。

\textbf{优化依据。} 二分前必须先定义谓词：某个位置是否已经满足边界条件，或者某个答案是否可行。数组二分找的是边界，答案二分找的是最小可行值或最大可行值。 本题的优化要抓住“当前值大于目标就左移，小于目标就下移。”这一点，把原来重复扫描或重复搜索的部分改成增量维护。优化以后，每次循环都应该只处理新增状态、删除过期状态或合并两个已经稳定的子答案，而不是重新审查全部输入。若题目条件发生变化，必须重新检查这个依据是否还成立。

\textbf{正确性不变量。} 建议固定左闭右开区间。循环中始终维护答案在 \code{[left, right)} 内，左侧一定不可行或小于目标，右侧外部已经被排除。 用这道题复盘时，可以把不变量写成三句话：已经处理的部分满足什么，尚未处理的部分还保留什么可能性，当前操作为什么不会丢掉最优答案或合法答案。证明二分不是证明循环会停，而是证明被排除的一半确实不含答案。答案空间二分还要证明可行性谓词单调。 证明不需要很形式化，但必须能排除反例；如果你说“感觉这样选最好”，就还没有完成算法说明。

\textbf{C++ 实现要点。} 中点写成 \code{left + (right - left) / 2}。判断平方、乘法、总量时避免溢出。用 \code{std::lower_bound} 能表达清楚时优先用它，但面试要能手写同等语义。本题还要特别注意：空矩阵、单行、单列、目标在角落、重复值。写代码时先把边界条件放在函数开头或循环不变量里，而不是用很多补丁式 if 修修补补。容器选择要和访问模式一致：需要顺序就别用无序表，需要极值就别完整排序，需要历史计数就别只保存布尔值。

\textbf{复杂度和边界。} 时间复杂度要按真实输入维度说明，不能把字符串长度、图边数、值域大小都混成一个 n。空间复杂度也要区分辅助结构和输出本身。测试时至少覆盖：最小输入、没有答案、答案在边界、重复元素或重复状态、以及会让暴力解最慢的规模输入。对于这题，边界复盘重点是：空矩阵、单行、单列、目标在角落、重复值。

\textbf{扩展和实验。} 从左上角无法一次排除一行或一列。实验准备长度为零、一、二的数组，目标小于全部元素、大于全部元素、重复元素边界、旋转数组、不可行答案和刚好可行答案。实验报告建议记录三件事：暴力版本在什么规模开始变慢，优化版本减少了哪类重复工作，哪一个边界用例最容易打破错误实现。这样一张题卡才算从“看过题解”变成“能够迁移”。

\subsection{704 二分查找}

\textbf{题目模型。} 基础题的重点是固定区间语义。这题归入“二分查找与答案空间”主线，因为它的核心不是某个语法技巧，而是如何把输入状态组织成可维护的结构。读题时先把对象、关系、约束和输出目标分开：对象决定容器，关系决定边或区间，约束决定能否单调推进，输出目标决定保存最值、计数、路径还是判定结果。左闭右开写法能统一 lower bound 和存在性检查。

\textbf{暴力解和瓶颈。} 暴力做法通常线性扫描数组或线性枚举答案。它的问题不是不会找到答案，而是没有利用有序性或可行性的单调边界。 对于本题，先写暴力解的价值在于看清状态空间：哪些候选会被枚举，哪些信息会被重复计算，哪些检查其实只依赖少量历史状态。慢点在于每次只排除一个候选。二分能成立，是因为一次比较可以排除一半下标，或者一次可行性检查可以排除一半答案范围。 如果面试中直接跳到优化而说不清暴力慢在哪里，后续复杂度分析就会显得像背模板。

\textbf{优化依据。} 二分前必须先定义谓词：某个位置是否已经满足边界条件，或者某个答案是否可行。数组二分找的是边界，答案二分找的是最小可行值或最大可行值。 本题的优化要抓住“左闭右开写法能统一 lower bound 和存在性检查。”这一点，把原来重复扫描或重复搜索的部分改成增量维护。优化以后，每次循环都应该只处理新增状态、删除过期状态或合并两个已经稳定的子答案，而不是重新审查全部输入。若题目条件发生变化，必须重新检查这个依据是否还成立。

\textbf{正确性不变量。} 建议固定左闭右开区间。循环中始终维护答案在 \code{[left, right)} 内，左侧一定不可行或小于目标，右侧外部已经被排除。 用这道题复盘时，可以把不变量写成三句话：已经处理的部分满足什么，尚未处理的部分还保留什么可能性，当前操作为什么不会丢掉最优答案或合法答案。证明二分不是证明循环会停，而是证明被排除的一半确实不含答案。答案空间二分还要证明可行性谓词单调。 证明不需要很形式化，但必须能排除反例；如果你说“感觉这样选最好”，就还没有完成算法说明。

\textbf{C++ 实现要点。} 中点写成 \code{left + (right - left) / 2}。判断平方、乘法、总量时避免溢出。用 \code{std::lower_bound} 能表达清楚时优先用它，但面试要能手写同等语义。本题还要特别注意：空数组、长度一、目标在边界、目标不存在。写代码时先把边界条件放在函数开头或循环不变量里，而不是用很多补丁式 if 修修补补。容器选择要和访问模式一致：需要顺序就别用无序表，需要极值就别完整排序，需要历史计数就别只保存布尔值。

\textbf{复杂度和边界。} 时间复杂度要按真实输入维度说明，不能把字符串长度、图边数、值域大小都混成一个 n。空间复杂度也要区分辅助结构和输出本身。测试时至少覆盖：最小输入、没有答案、答案在边界、重复元素或重复状态、以及会让暴力解最慢的规模输入。对于这题，边界复盘重点是：空数组、长度一、目标在边界、目标不存在。

\textbf{扩展和实验。} 所有复杂二分都应该能退回到这道题的边界不变量。实验准备长度为零、一、二的数组，目标小于全部元素、大于全部元素、重复元素边界、旋转数组、不可行答案和刚好可行答案。实验报告建议记录三件事：暴力版本在什么规模开始变慢，优化版本减少了哪类重复工作，哪一个边界用例最容易打破错误实现。这样一张题卡才算从“看过题解”变成“能够迁移”。

\subsection{875 爱吃香蕉的珂珂}

\textbf{题目模型。} 速度越大，总耗时越小，答案具有单调可行性。这题归入“二分查找与答案空间”主线，因为它的核心不是某个语法技巧，而是如何把输入状态组织成可维护的结构。读题时先把对象、关系、约束和输出目标分开：对象决定容器，关系决定边或区间，约束决定能否单调推进，输出目标决定保存最值、计数、路径还是判定结果。二分速度，检查函数用向上取整累加小时数。

\textbf{暴力解和瓶颈。} 暴力做法通常线性扫描数组或线性枚举答案。它的问题不是不会找到答案，而是没有利用有序性或可行性的单调边界。 对于本题，先写暴力解的价值在于看清状态空间：哪些候选会被枚举，哪些信息会被重复计算，哪些检查其实只依赖少量历史状态。慢点在于每次只排除一个候选。二分能成立，是因为一次比较可以排除一半下标，或者一次可行性检查可以排除一半答案范围。 如果面试中直接跳到优化而说不清暴力慢在哪里，后续复杂度分析就会显得像背模板。

\textbf{优化依据。} 二分前必须先定义谓词：某个位置是否已经满足边界条件，或者某个答案是否可行。数组二分找的是边界，答案二分找的是最小可行值或最大可行值。 本题的优化要抓住“二分速度，检查函数用向上取整累加小时数。”这一点，把原来重复扫描或重复搜索的部分改成增量维护。优化以后，每次循环都应该只处理新增状态、删除过期状态或合并两个已经稳定的子答案，而不是重新审查全部输入。若题目条件发生变化，必须重新检查这个依据是否还成立。

\textbf{正确性不变量。} 建议固定左闭右开区间。循环中始终维护答案在 \code{[left, right)} 内，左侧一定不可行或小于目标，右侧外部已经被排除。 用这道题复盘时，可以把不变量写成三句话：已经处理的部分满足什么，尚未处理的部分还保留什么可能性，当前操作为什么不会丢掉最优答案或合法答案。证明二分不是证明循环会停，而是证明被排除的一半确实不含答案。答案空间二分还要证明可行性谓词单调。 证明不需要很形式化，但必须能排除反例；如果你说“感觉这样选最好”，就还没有完成算法说明。

\textbf{C++ 实现要点。} 中点写成 \code{left + (right - left) / 2}。判断平方、乘法、总量时避免溢出。用 \code{std::lower_bound} 能表达清楚时优先用它，但面试要能手写同等语义。本题还要特别注意：速度下界、堆很大、h 等于堆数、乘法加法溢出。写代码时先把边界条件放在函数开头或循环不变量里，而不是用很多补丁式 if 修修补补。容器选择要和访问模式一致：需要顺序就别用无序表，需要极值就别完整排序，需要历史计数就别只保存布尔值。

\textbf{复杂度和边界。} 时间复杂度要按真实输入维度说明，不能把字符串长度、图边数、值域大小都混成一个 n。空间复杂度也要区分辅助结构和输出本身。测试时至少覆盖：最小输入、没有答案、答案在边界、重复元素或重复状态、以及会让暴力解最慢的规模输入。对于这题，边界复盘重点是：速度下界、堆很大、h 等于堆数、乘法加法溢出。

\textbf{扩展和实验。} 船运包裹和分割数组都是同类最小可行上限。实验准备长度为零、一、二的数组，目标小于全部元素、大于全部元素、重复元素边界、旋转数组、不可行答案和刚好可行答案。实验报告建议记录三件事：暴力版本在什么规模开始变慢，优化版本减少了哪类重复工作，哪一个边界用例最容易打破错误实现。这样一张题卡才算从“看过题解”变成“能够迁移”。

\subsection{1011 在 D 天内送达包裹}

\textbf{题目模型。} 容量越大，需要天数越少，存在最小可行容量。这题归入“二分查找与答案空间”主线，因为它的核心不是某个语法技巧，而是如何把输入状态组织成可维护的结构。读题时先把对象、关系、约束和输出目标分开：对象决定容器，关系决定边或区间，约束决定能否单调推进，输出目标决定保存最值、计数、路径还是判定结果。下界是最大包裹，上界是总重量，检查函数按顺序贪心装船。

\textbf{暴力解和瓶颈。} 暴力做法通常线性扫描数组或线性枚举答案。它的问题不是不会找到答案，而是没有利用有序性或可行性的单调边界。 对于本题，先写暴力解的价值在于看清状态空间：哪些候选会被枚举，哪些信息会被重复计算，哪些检查其实只依赖少量历史状态。慢点在于每次只排除一个候选。二分能成立，是因为一次比较可以排除一半下标，或者一次可行性检查可以排除一半答案范围。 如果面试中直接跳到优化而说不清暴力慢在哪里，后续复杂度分析就会显得像背模板。

\textbf{优化依据。} 二分前必须先定义谓词：某个位置是否已经满足边界条件，或者某个答案是否可行。数组二分找的是边界，答案二分找的是最小可行值或最大可行值。 本题的优化要抓住“下界是最大包裹，上界是总重量，检查函数按顺序贪心装船。”这一点，把原来重复扫描或重复搜索的部分改成增量维护。优化以后，每次循环都应该只处理新增状态、删除过期状态或合并两个已经稳定的子答案，而不是重新审查全部输入。若题目条件发生变化，必须重新检查这个依据是否还成立。

\textbf{正确性不变量。} 建议固定左闭右开区间。循环中始终维护答案在 \code{[left, right)} 内，左侧一定不可行或小于目标，右侧外部已经被排除。 用这道题复盘时，可以把不变量写成三句话：已经处理的部分满足什么，尚未处理的部分还保留什么可能性，当前操作为什么不会丢掉最优答案或合法答案。证明二分不是证明循环会停，而是证明被排除的一半确实不含答案。答案空间二分还要证明可行性谓词单调。 证明不需要很形式化，但必须能排除反例；如果你说“感觉这样选最好”，就还没有完成算法说明。

\textbf{C++ 实现要点。} 中点写成 \code{left + (right - left) / 2}。判断平方、乘法、总量时避免溢出。用 \code{std::lower_bound} 能表达清楚时优先用它，但面试要能手写同等语义。本题还要特别注意：D 为一、D 等于包裹数、单个包裹超过猜测容量。写代码时先把边界条件放在函数开头或循环不变量里，而不是用很多补丁式 if 修修补补。容器选择要和访问模式一致：需要顺序就别用无序表，需要极值就别完整排序，需要历史计数就别只保存布尔值。

\textbf{复杂度和边界。} 时间复杂度要按真实输入维度说明，不能把字符串长度、图边数、值域大小都混成一个 n。空间复杂度也要区分辅助结构和输出本身。测试时至少覆盖：最小输入、没有答案、答案在边界、重复元素或重复状态、以及会让暴力解最慢的规模输入。对于这题，边界复盘重点是：D 为一、D 等于包裹数、单个包裹超过猜测容量。

\textbf{扩展和实验。} 它和分割数组最大值都是连续分段最大和模型。实验准备长度为零、一、二的数组，目标小于全部元素、大于全部元素、重复元素边界、旋转数组、不可行答案和刚好可行答案。实验报告建议记录三件事：暴力版本在什么规模开始变慢，优化版本减少了哪类重复工作，哪一个边界用例最容易打破错误实现。这样一张题卡才算从“看过题解”变成“能够迁移”。

\subsection{栈、单调栈与表达式}
下面的题卡都按同一条链路展开：题目模型、暴力解、瓶颈、优化依据、不变量、C++ 实现、边界和扩展。格式统一是为了训练面试表达，而每张卡的关键观察和失效条件都要单独复盘。

\subsection{42 接雨水}

\textbf{题目模型。} 每个位置的水由左侧最高和右侧最高的较小者决定。这题归入“栈、单调栈与表达式”主线，因为它的核心不是某个语法技巧，而是如何把输入状态组织成可维护的结构。读题时先把对象、关系、约束和输出目标分开：对象决定容器，关系决定边或区间，约束决定能否单调推进，输出目标决定保存最值、计数、路径还是判定结果。前后缀、双指针、单调栈都是不同状态维护方式；要能说明各自结算对象。

\textbf{暴力解和瓶颈。} 暴力做法对每个位置向左或向右寻找第一个满足条件的元素，或者反复删除已经匹配的局部结构。这会让同一个元素被比较很多次。 对于本题，先写暴力解的价值在于看清状态空间：哪些候选会被枚举，哪些信息会被重复计算，哪些检查其实只依赖少量历史状态。瓶颈是未解决元素没有被组织起来。单调栈保存的是仍在等待某个未来元素解决的候选；普通栈保存的是最近打开但尚未关闭的结构。 如果面试中直接跳到优化而说不清暴力慢在哪里，后续复杂度分析就会显得像背模板。

\textbf{优化依据。} 当新元素到来时，它会一次性解决栈顶一串被它支配的旧元素。被弹出的元素以后再也不会参与比较，因此总体线性。 本题的优化要抓住“前后缀、双指针、单调栈都是不同状态维护方式；要能说明各自结算对象。”这一点，把原来重复扫描或重复搜索的部分改成增量维护。优化以后，每次循环都应该只处理新增状态、删除过期状态或合并两个已经稳定的子答案，而不是重新审查全部输入。若题目条件发生变化，必须重新检查这个依据是否还成立。

\textbf{正确性不变量。} 栈内下标对应的值要保持单调，或者栈顶必须是最近未匹配对象。弹出时要说清楚当前元素充当右边界、下一个栈顶充当左边界，还是当前运算符触发计算。 用这道题复盘时，可以把不变量写成三句话：已经处理的部分满足什么，尚未处理的部分还保留什么可能性，当前操作为什么不会丢掉最优答案或合法答案。每个元素最多入栈一次、出栈一次，所以时间线性。正确性来自弹出时机：一旦遇到破坏单调性的元素，栈顶元素的答案已经确定。 证明不需要很形式化，但必须能排除反例；如果你说“感觉这样选最好”，就还没有完成算法说明。

\textbf{C++ 实现要点。} 栈里通常存下标而不是值，因为要计算距离、宽度和过期。可以用 \code{std::vector<int>} 当栈，便于查看顶部和减少适配器限制。本题还要特别注意：单调递增、单调递减、平台高度、多个凹槽、数组长度不足三。写代码时先把边界条件放在函数开头或循环不变量里，而不是用很多补丁式 if 修修补补。容器选择要和访问模式一致：需要顺序就别用无序表，需要极值就别完整排序，需要历史计数就别只保存布尔值。

\textbf{复杂度和边界。} 时间复杂度要按真实输入维度说明，不能把字符串长度、图边数、值域大小都混成一个 n。空间复杂度也要区分辅助结构和输出本身。测试时至少覆盖：最小输入、没有答案、答案在边界、重复元素或重复状态、以及会让暴力解最慢的规模输入。对于这题，边界复盘重点是：单调递增、单调递减、平台高度、多个凹槽、数组长度不足三。

\textbf{扩展和实验。} 若柱子变成二维网格，需要最小堆从边界向内扩展。实验应包含严格递增、严格递减、相等元素、哨兵边界、空栈弹出、减法除法操作数顺序等情况。实验报告建议记录三件事：暴力版本在什么规模开始变慢，优化版本减少了哪类重复工作，哪一个边界用例最容易打破错误实现。这样一张题卡才算从“看过题解”变成“能够迁移”。

\subsection{84 柱状图中最大的矩形}

\textbf{题目模型。} 每根柱子作为最矮高度时，需要知道左右第一个更矮位置。这题归入“栈、单调栈与表达式”主线，因为它的核心不是某个语法技巧，而是如何把输入状态组织成可维护的结构。读题时先把对象、关系、约束和输出目标分开：对象决定容器，关系决定边或区间，约束决定能否单调推进，输出目标决定保存最值、计数、路径还是判定结果。单调递增栈在遇到更矮柱时结算被弹出柱子的最大宽度。

\textbf{暴力解和瓶颈。} 暴力做法对每个位置向左或向右寻找第一个满足条件的元素，或者反复删除已经匹配的局部结构。这会让同一个元素被比较很多次。 对于本题，先写暴力解的价值在于看清状态空间：哪些候选会被枚举，哪些信息会被重复计算，哪些检查其实只依赖少量历史状态。瓶颈是未解决元素没有被组织起来。单调栈保存的是仍在等待某个未来元素解决的候选；普通栈保存的是最近打开但尚未关闭的结构。 如果面试中直接跳到优化而说不清暴力慢在哪里，后续复杂度分析就会显得像背模板。

\textbf{优化依据。} 当新元素到来时，它会一次性解决栈顶一串被它支配的旧元素。被弹出的元素以后再也不会参与比较，因此总体线性。 本题的优化要抓住“单调递增栈在遇到更矮柱时结算被弹出柱子的最大宽度。”这一点，把原来重复扫描或重复搜索的部分改成增量维护。优化以后，每次循环都应该只处理新增状态、删除过期状态或合并两个已经稳定的子答案，而不是重新审查全部输入。若题目条件发生变化，必须重新检查这个依据是否还成立。

\textbf{正确性不变量。} 栈内下标对应的值要保持单调，或者栈顶必须是最近未匹配对象。弹出时要说清楚当前元素充当右边界、下一个栈顶充当左边界，还是当前运算符触发计算。 用这道题复盘时，可以把不变量写成三句话：已经处理的部分满足什么，尚未处理的部分还保留什么可能性，当前操作为什么不会丢掉最优答案或合法答案。每个元素最多入栈一次、出栈一次，所以时间线性。正确性来自弹出时机：一旦遇到破坏单调性的元素，栈顶元素的答案已经确定。 证明不需要很形式化，但必须能排除反例；如果你说“感觉这样选最好”，就还没有完成算法说明。

\textbf{C++ 实现要点。} 栈里通常存下标而不是值，因为要计算距离、宽度和过期。可以用 \code{std::vector<int>} 当栈，便于查看顶部和减少适配器限制。本题还要特别注意：递增、递减、相等高度、哨兵零、宽度计算。写代码时先把边界条件放在函数开头或循环不变量里，而不是用很多补丁式 if 修修补补。容器选择要和访问模式一致：需要顺序就别用无序表，需要极值就别完整排序，需要历史计数就别只保存布尔值。

\textbf{复杂度和边界。} 时间复杂度要按真实输入维度说明，不能把字符串长度、图边数、值域大小都混成一个 n。空间复杂度也要区分辅助结构和输出本身。测试时至少覆盖：最小输入、没有答案、答案在边界、重复元素或重复状态、以及会让暴力解最慢的规模输入。对于这题，边界复盘重点是：递增、递减、相等高度、哨兵零、宽度计算。

\textbf{扩展和实验。} 最大矩形可以作为二维矩阵最大矩形的行压缩子问题。实验应包含严格递增、严格递减、相等元素、哨兵边界、空栈弹出、减法除法操作数顺序等情况。实验报告建议记录三件事：暴力版本在什么规模开始变慢，优化版本减少了哪类重复工作，哪一个边界用例最容易打破错误实现。这样一张题卡才算从“看过题解”变成“能够迁移”。

\subsection{150 逆波兰表达式求值}

\textbf{题目模型。} 后缀表达式把优先级和括号都编码进 token 顺序。这题归入“栈、单调栈与表达式”主线，因为它的核心不是某个语法技巧，而是如何把输入状态组织成可维护的结构。读题时先把对象、关系、约束和输出目标分开：对象决定容器，关系决定边或区间，约束决定能否单调推进，输出目标决定保存最值、计数、路径还是判定结果。遇到数字入栈，遇到运算符弹出右操作数和左操作数计算。

\textbf{暴力解和瓶颈。} 暴力做法对每个位置向左或向右寻找第一个满足条件的元素，或者反复删除已经匹配的局部结构。这会让同一个元素被比较很多次。 对于本题，先写暴力解的价值在于看清状态空间：哪些候选会被枚举，哪些信息会被重复计算，哪些检查其实只依赖少量历史状态。瓶颈是未解决元素没有被组织起来。单调栈保存的是仍在等待某个未来元素解决的候选；普通栈保存的是最近打开但尚未关闭的结构。 如果面试中直接跳到优化而说不清暴力慢在哪里，后续复杂度分析就会显得像背模板。

\textbf{优化依据。} 当新元素到来时，它会一次性解决栈顶一串被它支配的旧元素。被弹出的元素以后再也不会参与比较，因此总体线性。 本题的优化要抓住“遇到数字入栈，遇到运算符弹出右操作数和左操作数计算。”这一点，把原来重复扫描或重复搜索的部分改成增量维护。优化以后，每次循环都应该只处理新增状态、删除过期状态或合并两个已经稳定的子答案，而不是重新审查全部输入。若题目条件发生变化，必须重新检查这个依据是否还成立。

\textbf{正确性不变量。} 栈内下标对应的值要保持单调，或者栈顶必须是最近未匹配对象。弹出时要说清楚当前元素充当右边界、下一个栈顶充当左边界，还是当前运算符触发计算。 用这道题复盘时，可以把不变量写成三句话：已经处理的部分满足什么，尚未处理的部分还保留什么可能性，当前操作为什么不会丢掉最优答案或合法答案。每个元素最多入栈一次、出栈一次，所以时间线性。正确性来自弹出时机：一旦遇到破坏单调性的元素，栈顶元素的答案已经确定。 证明不需要很形式化，但必须能排除反例；如果你说“感觉这样选最好”，就还没有完成算法说明。

\textbf{C++ 实现要点。} 栈里通常存下标而不是值，因为要计算距离、宽度和过期。可以用 \code{std::vector<int>} 当栈，便于查看顶部和减少适配器限制。本题还要特别注意：负数 token、多位数、除法向零截断、减法顺序。写代码时先把边界条件放在函数开头或循环不变量里，而不是用很多补丁式 if 修修补补。容器选择要和访问模式一致：需要顺序就别用无序表，需要极值就别完整排序，需要历史计数就别只保存布尔值。

\textbf{复杂度和边界。} 时间复杂度要按真实输入维度说明，不能把字符串长度、图边数、值域大小都混成一个 n。空间复杂度也要区分辅助结构和输出本身。测试时至少覆盖：最小输入、没有答案、答案在边界、重复元素或重复状态、以及会让暴力解最慢的规模输入。对于这题，边界复盘重点是：负数 token、多位数、除法向零截断、减法顺序。

\textbf{扩展和实验。} 中缀表达式求值需要另一个运算符栈处理优先级。实验应包含严格递增、严格递减、相等元素、哨兵边界、空栈弹出、减法除法操作数顺序等情况。实验报告建议记录三件事：暴力版本在什么规模开始变慢，优化版本减少了哪类重复工作，哪一个边界用例最容易打破错误实现。这样一张题卡才算从“看过题解”变成“能够迁移”。

\subsection{155 最小栈}

\textbf{题目模型。} 普通栈只能看到栈顶，不能常数时间知道历史最小值。这题归入“栈、单调栈与表达式”主线，因为它的核心不是某个语法技巧，而是如何把输入状态组织成可维护的结构。读题时先把对象、关系、约束和输出目标分开：对象决定容器，关系决定边或区间，约束决定能否单调推进，输出目标决定保存最值、计数、路径还是判定结果。辅助栈同步保存每个深度对应的最小值。

\textbf{暴力解和瓶颈。} 暴力做法对每个位置向左或向右寻找第一个满足条件的元素，或者反复删除已经匹配的局部结构。这会让同一个元素被比较很多次。 对于本题，先写暴力解的价值在于看清状态空间：哪些候选会被枚举，哪些信息会被重复计算，哪些检查其实只依赖少量历史状态。瓶颈是未解决元素没有被组织起来。单调栈保存的是仍在等待某个未来元素解决的候选；普通栈保存的是最近打开但尚未关闭的结构。 如果面试中直接跳到优化而说不清暴力慢在哪里，后续复杂度分析就会显得像背模板。

\textbf{优化依据。} 当新元素到来时，它会一次性解决栈顶一串被它支配的旧元素。被弹出的元素以后再也不会参与比较，因此总体线性。 本题的优化要抓住“辅助栈同步保存每个深度对应的最小值。”这一点，把原来重复扫描或重复搜索的部分改成增量维护。优化以后，每次循环都应该只处理新增状态、删除过期状态或合并两个已经稳定的子答案，而不是重新审查全部输入。若题目条件发生变化，必须重新检查这个依据是否还成立。

\textbf{正确性不变量。} 栈内下标对应的值要保持单调，或者栈顶必须是最近未匹配对象。弹出时要说清楚当前元素充当右边界、下一个栈顶充当左边界，还是当前运算符触发计算。 用这道题复盘时，可以把不变量写成三句话：已经处理的部分满足什么，尚未处理的部分还保留什么可能性，当前操作为什么不会丢掉最优答案或合法答案。每个元素最多入栈一次、出栈一次，所以时间线性。正确性来自弹出时机：一旦遇到破坏单调性的元素，栈顶元素的答案已经确定。 证明不需要很形式化，但必须能排除反例；如果你说“感觉这样选最好”，就还没有完成算法说明。

\textbf{C++ 实现要点。} 栈里通常存下标而不是值，因为要计算距离、宽度和过期。可以用 \code{std::vector<int>} 当栈，便于查看顶部和减少适配器限制。本题还要特别注意：重复最小值、弹出最小值、空栈操作约定。写代码时先把边界条件放在函数开头或循环不变量里，而不是用很多补丁式 if 修修补补。容器选择要和访问模式一致：需要顺序就别用无序表，需要极值就别完整排序，需要历史计数就别只保存布尔值。

\textbf{复杂度和边界。} 时间复杂度要按真实输入维度说明，不能把字符串长度、图边数、值域大小都混成一个 n。空间复杂度也要区分辅助结构和输出本身。测试时至少覆盖：最小输入、没有答案、答案在边界、重复元素或重复状态、以及会让暴力解最慢的规模输入。对于这题，边界复盘重点是：重复最小值、弹出最小值、空栈操作约定。

\textbf{扩展和实验。} 也可保存差值压缩空间，但可读性和溢出处理更复杂。实验应包含严格递增、严格递减、相等元素、哨兵边界、空栈弹出、减法除法操作数顺序等情况。实验报告建议记录三件事：暴力版本在什么规模开始变慢，优化版本减少了哪类重复工作，哪一个边界用例最容易打破错误实现。这样一张题卡才算从“看过题解”变成“能够迁移”。

\subsection{232 用栈实现队列}

\textbf{题目模型。} 两个栈分别负责输入顺序和输出顺序。这题归入“栈、单调栈与表达式”主线，因为它的核心不是某个语法技巧，而是如何把输入状态组织成可维护的结构。读题时先把对象、关系、约束和输出目标分开：对象决定容器，关系决定边或区间，约束决定能否单调推进，输出目标决定保存最值、计数、路径还是判定结果。只有输出栈为空时，才把输入栈整体倒过去，保证均摊常数。

\textbf{暴力解和瓶颈。} 暴力做法对每个位置向左或向右寻找第一个满足条件的元素，或者反复删除已经匹配的局部结构。这会让同一个元素被比较很多次。 对于本题，先写暴力解的价值在于看清状态空间：哪些候选会被枚举，哪些信息会被重复计算，哪些检查其实只依赖少量历史状态。瓶颈是未解决元素没有被组织起来。单调栈保存的是仍在等待某个未来元素解决的候选；普通栈保存的是最近打开但尚未关闭的结构。 如果面试中直接跳到优化而说不清暴力慢在哪里，后续复杂度分析就会显得像背模板。

\textbf{优化依据。} 当新元素到来时，它会一次性解决栈顶一串被它支配的旧元素。被弹出的元素以后再也不会参与比较，因此总体线性。 本题的优化要抓住“只有输出栈为空时，才把输入栈整体倒过去，保证均摊常数。”这一点，把原来重复扫描或重复搜索的部分改成增量维护。优化以后，每次循环都应该只处理新增状态、删除过期状态或合并两个已经稳定的子答案，而不是重新审查全部输入。若题目条件发生变化，必须重新检查这个依据是否还成立。

\textbf{正确性不变量。} 栈内下标对应的值要保持单调，或者栈顶必须是最近未匹配对象。弹出时要说清楚当前元素充当右边界、下一个栈顶充当左边界，还是当前运算符触发计算。 用这道题复盘时，可以把不变量写成三句话：已经处理的部分满足什么，尚未处理的部分还保留什么可能性，当前操作为什么不会丢掉最优答案或合法答案。每个元素最多入栈一次、出栈一次，所以时间线性。正确性来自弹出时机：一旦遇到破坏单调性的元素，栈顶元素的答案已经确定。 证明不需要很形式化，但必须能排除反例；如果你说“感觉这样选最好”，就还没有完成算法说明。

\textbf{C++ 实现要点。} 栈里通常存下标而不是值，因为要计算距离、宽度和过期。可以用 \code{std::vector<int>} 当栈，便于查看顶部和减少适配器限制。本题还要特别注意：连续 peek、空队列、交替 push pop。写代码时先把边界条件放在函数开头或循环不变量里，而不是用很多补丁式 if 修修补补。容器选择要和访问模式一致：需要顺序就别用无序表，需要极值就别完整排序，需要历史计数就别只保存布尔值。

\textbf{复杂度和边界。} 时间复杂度要按真实输入维度说明，不能把字符串长度、图边数、值域大小都混成一个 n。空间复杂度也要区分辅助结构和输出本身。测试时至少覆盖：最小输入、没有答案、答案在边界、重复元素或重复状态、以及会让暴力解最慢的规模输入。对于这题，边界复盘重点是：连续 peek、空队列、交替 push pop。

\textbf{扩展和实验。} 用队列实现栈则要通过轮转保持最近元素在队首。实验应包含严格递增、严格递减、相等元素、哨兵边界、空栈弹出、减法除法操作数顺序等情况。实验报告建议记录三件事：暴力版本在什么规模开始变慢，优化版本减少了哪类重复工作，哪一个边界用例最容易打破错误实现。这样一张题卡才算从“看过题解”变成“能够迁移”。

\subsection{239 滑动窗口最大值}

\textbf{题目模型。} 每个窗口最大值来自仍未过期且没有被新元素支配的候选。这题归入“栈、单调栈与表达式”主线，因为它的核心不是某个语法技巧，而是如何把输入状态组织成可维护的结构。读题时先把对象、关系、约束和输出目标分开：对象决定容器，关系决定边或区间，约束决定能否单调推进，输出目标决定保存最值、计数、路径还是判定结果。单调队列保存下标，队尾小于等于新值的候选被弹出。

\textbf{暴力解和瓶颈。} 暴力做法对每个位置向左或向右寻找第一个满足条件的元素，或者反复删除已经匹配的局部结构。这会让同一个元素被比较很多次。 对于本题，先写暴力解的价值在于看清状态空间：哪些候选会被枚举，哪些信息会被重复计算，哪些检查其实只依赖少量历史状态。瓶颈是未解决元素没有被组织起来。单调栈保存的是仍在等待某个未来元素解决的候选；普通栈保存的是最近打开但尚未关闭的结构。 如果面试中直接跳到优化而说不清暴力慢在哪里，后续复杂度分析就会显得像背模板。

\textbf{优化依据。} 当新元素到来时，它会一次性解决栈顶一串被它支配的旧元素。被弹出的元素以后再也不会参与比较，因此总体线性。 本题的优化要抓住“单调队列保存下标，队尾小于等于新值的候选被弹出。”这一点，把原来重复扫描或重复搜索的部分改成增量维护。优化以后，每次循环都应该只处理新增状态、删除过期状态或合并两个已经稳定的子答案，而不是重新审查全部输入。若题目条件发生变化，必须重新检查这个依据是否还成立。

\textbf{正确性不变量。} 栈内下标对应的值要保持单调，或者栈顶必须是最近未匹配对象。弹出时要说清楚当前元素充当右边界、下一个栈顶充当左边界，还是当前运算符触发计算。 用这道题复盘时，可以把不变量写成三句话：已经处理的部分满足什么，尚未处理的部分还保留什么可能性，当前操作为什么不会丢掉最优答案或合法答案。每个元素最多入栈一次、出栈一次，所以时间线性。正确性来自弹出时机：一旦遇到破坏单调性的元素，栈顶元素的答案已经确定。 证明不需要很形式化，但必须能排除反例；如果你说“感觉这样选最好”，就还没有完成算法说明。

\textbf{C++ 实现要点。} 栈里通常存下标而不是值，因为要计算距离、宽度和过期。可以用 \code{std::vector<int>} 当栈，便于查看顶部和减少适配器限制。本题还要特别注意：k 为一、k 等于数组长度、重复最大值、队首过期。写代码时先把边界条件放在函数开头或循环不变量里，而不是用很多补丁式 if 修修补补。容器选择要和访问模式一致：需要顺序就别用无序表，需要极值就别完整排序，需要历史计数就别只保存布尔值。

\textbf{复杂度和边界。} 时间复杂度要按真实输入维度说明，不能把字符串长度、图边数、值域大小都混成一个 n。空间复杂度也要区分辅助结构和输出本身。测试时至少覆盖：最小输入、没有答案、答案在边界、重复元素或重复状态、以及会让暴力解最慢的规模输入。对于这题，边界复盘重点是：k 为一、k 等于数组长度、重复最大值、队首过期。

\textbf{扩展和实验。} 受限子序列和使用同样队列维护最近 k 个 DP 最大值。实验应包含严格递增、严格递减、相等元素、哨兵边界、空栈弹出、减法除法操作数顺序等情况。实验报告建议记录三件事：暴力版本在什么规模开始变慢，优化版本减少了哪类重复工作，哪一个边界用例最容易打破错误实现。这样一张题卡才算从“看过题解”变成“能够迁移”。

\subsection{394 字符串解码}

\textbf{题目模型。} 嵌套结构需要保存上一层字符串和重复次数。这题归入“栈、单调栈与表达式”主线，因为它的核心不是某个语法技巧，而是如何把输入状态组织成可维护的结构。读题时先把对象、关系、约束和输出目标分开：对象决定容器，关系决定边或区间，约束决定能否单调推进，输出目标决定保存最值、计数、路径还是判定结果。遇到左括号入栈当前状态，右括号弹出并拼接展开。

\textbf{暴力解和瓶颈。} 暴力做法对每个位置向左或向右寻找第一个满足条件的元素，或者反复删除已经匹配的局部结构。这会让同一个元素被比较很多次。 对于本题，先写暴力解的价值在于看清状态空间：哪些候选会被枚举，哪些信息会被重复计算，哪些检查其实只依赖少量历史状态。瓶颈是未解决元素没有被组织起来。单调栈保存的是仍在等待某个未来元素解决的候选；普通栈保存的是最近打开但尚未关闭的结构。 如果面试中直接跳到优化而说不清暴力慢在哪里，后续复杂度分析就会显得像背模板。

\textbf{优化依据。} 当新元素到来时，它会一次性解决栈顶一串被它支配的旧元素。被弹出的元素以后再也不会参与比较，因此总体线性。 本题的优化要抓住“遇到左括号入栈当前状态，右括号弹出并拼接展开。”这一点，把原来重复扫描或重复搜索的部分改成增量维护。优化以后，每次循环都应该只处理新增状态、删除过期状态或合并两个已经稳定的子答案，而不是重新审查全部输入。若题目条件发生变化，必须重新检查这个依据是否还成立。

\textbf{正确性不变量。} 栈内下标对应的值要保持单调，或者栈顶必须是最近未匹配对象。弹出时要说清楚当前元素充当右边界、下一个栈顶充当左边界，还是当前运算符触发计算。 用这道题复盘时，可以把不变量写成三句话：已经处理的部分满足什么，尚未处理的部分还保留什么可能性，当前操作为什么不会丢掉最优答案或合法答案。每个元素最多入栈一次、出栈一次，所以时间线性。正确性来自弹出时机：一旦遇到破坏单调性的元素，栈顶元素的答案已经确定。 证明不需要很形式化，但必须能排除反例；如果你说“感觉这样选最好”，就还没有完成算法说明。

\textbf{C++ 实现要点。} 栈里通常存下标而不是值，因为要计算距离、宽度和过期。可以用 \code{std::vector<int>} 当栈，便于查看顶部和减少适配器限制。本题还要特别注意：多位数字、嵌套、空片段、数字后必须跟括号。写代码时先把边界条件放在函数开头或循环不变量里，而不是用很多补丁式 if 修修补补。容器选择要和访问模式一致：需要顺序就别用无序表，需要极值就别完整排序，需要历史计数就别只保存布尔值。

\textbf{复杂度和边界。} 时间复杂度要按真实输入维度说明，不能把字符串长度、图边数、值域大小都混成一个 n。空间复杂度也要区分辅助结构和输出本身。测试时至少覆盖：最小输入、没有答案、答案在边界、重复元素或重复状态、以及会让暴力解最慢的规模输入。对于这题，边界复盘重点是：多位数字、嵌套、空片段、数字后必须跟括号。

\textbf{扩展和实验。} 递归下降解析更通用，但迭代栈更符合工程栈控制。实验应包含严格递增、严格递减、相等元素、哨兵边界、空栈弹出、减法除法操作数顺序等情况。实验报告建议记录三件事：暴力版本在什么规模开始变慢，优化版本减少了哪类重复工作，哪一个边界用例最容易打破错误实现。这样一张题卡才算从“看过题解”变成“能够迁移”。

\subsection{496 下一个更大元素 I}

\textbf{题目模型。} 第二个数组提供每个值右侧第一个更大值的全局关系。这题归入“栈、单调栈与表达式”主线，因为它的核心不是某个语法技巧，而是如何把输入状态组织成可维护的结构。读题时先把对象、关系、约束和输出目标分开：对象决定容器，关系决定边或区间，约束决定能否单调推进，输出目标决定保存最值、计数、路径还是判定结果。单调栈预处理值到下一个更大值映射，再回答第一个数组查询。

\textbf{暴力解和瓶颈。} 暴力做法对每个位置向左或向右寻找第一个满足条件的元素，或者反复删除已经匹配的局部结构。这会让同一个元素被比较很多次。 对于本题，先写暴力解的价值在于看清状态空间：哪些候选会被枚举，哪些信息会被重复计算，哪些检查其实只依赖少量历史状态。瓶颈是未解决元素没有被组织起来。单调栈保存的是仍在等待某个未来元素解决的候选；普通栈保存的是最近打开但尚未关闭的结构。 如果面试中直接跳到优化而说不清暴力慢在哪里，后续复杂度分析就会显得像背模板。

\textbf{优化依据。} 当新元素到来时，它会一次性解决栈顶一串被它支配的旧元素。被弹出的元素以后再也不会参与比较，因此总体线性。 本题的优化要抓住“单调栈预处理值到下一个更大值映射，再回答第一个数组查询。”这一点，把原来重复扫描或重复搜索的部分改成增量维护。优化以后，每次循环都应该只处理新增状态、删除过期状态或合并两个已经稳定的子答案，而不是重新审查全部输入。若题目条件发生变化，必须重新检查这个依据是否还成立。

\textbf{正确性不变量。} 栈内下标对应的值要保持单调，或者栈顶必须是最近未匹配对象。弹出时要说清楚当前元素充当右边界、下一个栈顶充当左边界，还是当前运算符触发计算。 用这道题复盘时，可以把不变量写成三句话：已经处理的部分满足什么，尚未处理的部分还保留什么可能性，当前操作为什么不会丢掉最优答案或合法答案。每个元素最多入栈一次、出栈一次，所以时间线性。正确性来自弹出时机：一旦遇到破坏单调性的元素，栈顶元素的答案已经确定。 证明不需要很形式化，但必须能排除反例；如果你说“感觉这样选最好”，就还没有完成算法说明。

\textbf{C++ 实现要点。} 栈里通常存下标而不是值，因为要计算距离、宽度和过期。可以用 \code{std::vector<int>} 当栈，便于查看顶部和减少适配器限制。本题还要特别注意：不存在更大值、递减数组、值唯一假设。写代码时先把边界条件放在函数开头或循环不变量里，而不是用很多补丁式 if 修修补补。容器选择要和访问模式一致：需要顺序就别用无序表，需要极值就别完整排序，需要历史计数就别只保存布尔值。

\textbf{复杂度和边界。} 时间复杂度要按真实输入维度说明，不能把字符串长度、图边数、值域大小都混成一个 n。空间复杂度也要区分辅助结构和输出本身。测试时至少覆盖：最小输入、没有答案、答案在边界、重复元素或重复状态、以及会让暴力解最慢的规模输入。对于这题，边界复盘重点是：不存在更大值、递减数组、值唯一假设。

\textbf{扩展和实验。} 若值不唯一，应按下标而不是值建映射。实验应包含严格递增、严格递减、相等元素、哨兵边界、空栈弹出、减法除法操作数顺序等情况。实验报告建议记录三件事：暴力版本在什么规模开始变慢，优化版本减少了哪类重复工作，哪一个边界用例最容易打破错误实现。这样一张题卡才算从“看过题解”变成“能够迁移”。

\subsection{503 下一个更大元素 II}

\textbf{题目模型。} 环形数组可以通过遍历两倍下标模拟绕回。这题归入“栈、单调栈与表达式”主线，因为它的核心不是某个语法技巧，而是如何把输入状态组织成可维护的结构。读题时先把对象、关系、约束和输出目标分开：对象决定容器，关系决定边或区间，约束决定能否单调推进，输出目标决定保存最值、计数、路径还是判定结果。第一遍入栈，第二遍只帮助未解决下标找到答案。

\textbf{暴力解和瓶颈。} 暴力做法对每个位置向左或向右寻找第一个满足条件的元素，或者反复删除已经匹配的局部结构。这会让同一个元素被比较很多次。 对于本题，先写暴力解的价值在于看清状态空间：哪些候选会被枚举，哪些信息会被重复计算，哪些检查其实只依赖少量历史状态。瓶颈是未解决元素没有被组织起来。单调栈保存的是仍在等待某个未来元素解决的候选；普通栈保存的是最近打开但尚未关闭的结构。 如果面试中直接跳到优化而说不清暴力慢在哪里，后续复杂度分析就会显得像背模板。

\textbf{优化依据。} 当新元素到来时，它会一次性解决栈顶一串被它支配的旧元素。被弹出的元素以后再也不会参与比较，因此总体线性。 本题的优化要抓住“第一遍入栈，第二遍只帮助未解决下标找到答案。”这一点，把原来重复扫描或重复搜索的部分改成增量维护。优化以后，每次循环都应该只处理新增状态、删除过期状态或合并两个已经稳定的子答案，而不是重新审查全部输入。若题目条件发生变化，必须重新检查这个依据是否还成立。

\textbf{正确性不变量。} 栈内下标对应的值要保持单调，或者栈顶必须是最近未匹配对象。弹出时要说清楚当前元素充当右边界、下一个栈顶充当左边界，还是当前运算符触发计算。 用这道题复盘时，可以把不变量写成三句话：已经处理的部分满足什么，尚未处理的部分还保留什么可能性，当前操作为什么不会丢掉最优答案或合法答案。每个元素最多入栈一次、出栈一次，所以时间线性。正确性来自弹出时机：一旦遇到破坏单调性的元素，栈顶元素的答案已经确定。 证明不需要很形式化，但必须能排除反例；如果你说“感觉这样选最好”，就还没有完成算法说明。

\textbf{C++ 实现要点。} 栈里通常存下标而不是值，因为要计算距离、宽度和过期。可以用 \code{std::vector<int>} 当栈，便于查看顶部和减少适配器限制。本题还要特别注意：全递减、全相等、长度一、取模下标。写代码时先把边界条件放在函数开头或循环不变量里，而不是用很多补丁式 if 修修补补。容器选择要和访问模式一致：需要顺序就别用无序表，需要极值就别完整排序，需要历史计数就别只保存布尔值。

\textbf{复杂度和边界。} 时间复杂度要按真实输入维度说明，不能把字符串长度、图边数、值域大小都混成一个 n。空间复杂度也要区分辅助结构和输出本身。测试时至少覆盖：最小输入、没有答案、答案在边界、重复元素或重复状态、以及会让暴力解最慢的规模输入。对于这题，边界复盘重点是：全递减、全相等、长度一、取模下标。

\textbf{扩展和实验。} 不要在第二遍重复入栈，否则可能死循环或重复计算。实验应包含严格递增、严格递减、相等元素、哨兵边界、空栈弹出、减法除法操作数顺序等情况。实验报告建议记录三件事：暴力版本在什么规模开始变慢，优化版本减少了哪类重复工作，哪一个边界用例最容易打破错误实现。这样一张题卡才算从“看过题解”变成“能够迁移”。

\subsection{739 每日温度}

\textbf{题目模型。} 每一天等待右侧第一个更高温度。这题归入“栈、单调栈与表达式”主线，因为它的核心不是某个语法技巧，而是如何把输入状态组织成可维护的结构。读题时先把对象、关系、约束和输出目标分开：对象决定容器，关系决定边或区间，约束决定能否单调推进，输出目标决定保存最值、计数、路径还是判定结果。单调递减栈保存尚未找到答案的下标。

\textbf{暴力解和瓶颈。} 暴力做法对每个位置向左或向右寻找第一个满足条件的元素，或者反复删除已经匹配的局部结构。这会让同一个元素被比较很多次。 对于本题，先写暴力解的价值在于看清状态空间：哪些候选会被枚举，哪些信息会被重复计算，哪些检查其实只依赖少量历史状态。瓶颈是未解决元素没有被组织起来。单调栈保存的是仍在等待某个未来元素解决的候选；普通栈保存的是最近打开但尚未关闭的结构。 如果面试中直接跳到优化而说不清暴力慢在哪里，后续复杂度分析就会显得像背模板。

\textbf{优化依据。} 当新元素到来时，它会一次性解决栈顶一串被它支配的旧元素。被弹出的元素以后再也不会参与比较，因此总体线性。 本题的优化要抓住“单调递减栈保存尚未找到答案的下标。”这一点，把原来重复扫描或重复搜索的部分改成增量维护。优化以后，每次循环都应该只处理新增状态、删除过期状态或合并两个已经稳定的子答案，而不是重新审查全部输入。若题目条件发生变化，必须重新检查这个依据是否还成立。

\textbf{正确性不变量。} 栈内下标对应的值要保持单调，或者栈顶必须是最近未匹配对象。弹出时要说清楚当前元素充当右边界、下一个栈顶充当左边界，还是当前运算符触发计算。 用这道题复盘时，可以把不变量写成三句话：已经处理的部分满足什么，尚未处理的部分还保留什么可能性，当前操作为什么不会丢掉最优答案或合法答案。每个元素最多入栈一次、出栈一次，所以时间线性。正确性来自弹出时机：一旦遇到破坏单调性的元素，栈顶元素的答案已经确定。 证明不需要很形式化，但必须能排除反例；如果你说“感觉这样选最好”，就还没有完成算法说明。

\textbf{C++ 实现要点。} 栈里通常存下标而不是值，因为要计算距离、宽度和过期。可以用 \code{std::vector<int>} 当栈，便于查看顶部和减少适配器限制。本题还要特别注意：没有更高温、相等温度、递增、递减。写代码时先把边界条件放在函数开头或循环不变量里，而不是用很多补丁式 if 修修补补。容器选择要和访问模式一致：需要顺序就别用无序表，需要极值就别完整排序，需要历史计数就别只保存布尔值。

\textbf{复杂度和边界。} 时间复杂度要按真实输入维度说明，不能把字符串长度、图边数、值域大小都混成一个 n。空间复杂度也要区分辅助结构和输出本身。测试时至少覆盖：最小输入、没有答案、答案在边界、重复元素或重复状态、以及会让暴力解最慢的规模输入。对于这题，边界复盘重点是：没有更高温、相等温度、递增、递减。

\textbf{扩展和实验。} 下一个更大元素与它同型，只是输出值或距离不同。实验应包含严格递增、严格递减、相等元素、哨兵边界、空栈弹出、减法除法操作数顺序等情况。实验报告建议记录三件事：暴力版本在什么规模开始变慢，优化版本减少了哪类重复工作，哪一个边界用例最容易打破错误实现。这样一张题卡才算从“看过题解”变成“能够迁移”。

\subsection{位运算与状态压缩}
下面的题卡都按同一条链路展开：题目模型、暴力解、瓶颈、优化依据、不变量、C++ 实现、边界和扩展。格式统一是为了训练面试表达，而每张卡的关键观察和失效条件都要单独复盘。

\subsection{136 只出现一次的数字}

\textbf{题目模型。} 成对元素可以用异或抵消，剩下单个元素。这题归入“位运算与状态压缩”主线，因为它的核心不是某个语法技巧，而是如何把输入状态组织成可维护的结构。读题时先把对象、关系、约束和输出目标分开：对象决定容器，关系决定边或区间，约束决定能否单调推进，输出目标决定保存最值、计数、路径还是判定结果。异或满足交换结合，扫描顺序不影响答案。

\textbf{暴力解和瓶颈。} 暴力做法用集合、数组或递归保存状态，空间和比较成本较高。若状态本质是若干布尔选择，位掩码能把它压进整数。 对于本题，先写暴力解的价值在于看清状态空间：哪些候选会被枚举，哪些信息会被重复计算，哪些检查其实只依赖少量历史状态。慢点在于状态复制和集合比较。位操作让包含、加入、删除、枚举子集变成常数或接近常数的机器指令。 如果面试中直接跳到优化而说不清暴力慢在哪里，后续复杂度分析就会显得像背模板。

\textbf{优化依据。} 异或解决成对抵消，按位计数处理出现三次，掩码表示访问集合或可用集合，低位提取用于枚举候选。 本题的优化要抓住“异或满足交换结合，扫描顺序不影响答案。”这一点，把原来重复扫描或重复搜索的部分改成增量维护。优化以后，每次循环都应该只处理新增状态、删除过期状态或合并两个已经稳定的子答案，而不是重新审查全部输入。若题目条件发生变化，必须重新检查这个依据是否还成立。

\textbf{正确性不变量。} 掩码的每一位必须有固定含义。状态去重时，位置相同但掩码不同不能合并，因为未来能力不同。 用这道题复盘时，可以把不变量写成三句话：已经处理的部分满足什么，尚未处理的部分还保留什么可能性，当前操作为什么不会丢掉最优答案或合法答案。正确性来自二进制位独立性。异或按位满足交换结合和自反抵消；掩码 DP 则是对子集大小做归纳。 证明不需要很形式化，但必须能排除反例；如果你说“感觉这样选最好”，就还没有完成算法说明。

\textbf{C++ 实现要点。} 使用无符号整数能减少移位歧义。位数超过 31 时用 \code{long long} 或 \code{std::bitset}。表达式加括号，避免位运算优先级误读。本题还要特别注意：零、负数、唯一值在开头或结尾、所有其他值恰好两次。写代码时先把边界条件放在函数开头或循环不变量里，而不是用很多补丁式 if 修修补补。容器选择要和访问模式一致：需要顺序就别用无序表，需要极值就别完整排序，需要历史计数就别只保存布尔值。

\textbf{复杂度和边界。} 时间复杂度要按真实输入维度说明，不能把字符串长度、图边数、值域大小都混成一个 n。空间复杂度也要区分辅助结构和输出本身。测试时至少覆盖：最小输入、没有答案、答案在边界、重复元素或重复状态、以及会让暴力解最慢的规模输入。对于这题，边界复盘重点是：零、负数、唯一值在开头或结尾、所有其他值恰好两次。

\textbf{扩展和实验。} 若其他值出现三次，要改为按位计数取模。实验包含零、负数、最高位、重复数、全集掩码、空集合、只差一位的状态和资源耗尽后的不可达状态。实验报告建议记录三件事：暴力版本在什么规模开始变慢，优化版本减少了哪类重复工作，哪一个边界用例最容易打破错误实现。这样一张题卡才算从“看过题解”变成“能够迁移”。

\subsection{287 寻找重复数}

\textbf{题目模型。} 数组值域使下标到值形成函数图，重复数是环入口。这题归入“位运算与状态压缩”主线，因为它的核心不是某个语法技巧，而是如何把输入状态组织成可维护的结构。读题时先把对象、关系、约束和输出目标分开：对象决定容器，关系决定边或区间，约束决定能否单调推进，输出目标决定保存最值、计数、路径还是判定结果。快慢指针不修改数组且常数空间；值域二分用抽屉原理。

\textbf{暴力解和瓶颈。} 暴力做法用集合、数组或递归保存状态，空间和比较成本较高。若状态本质是若干布尔选择，位掩码能把它压进整数。 对于本题，先写暴力解的价值在于看清状态空间：哪些候选会被枚举，哪些信息会被重复计算，哪些检查其实只依赖少量历史状态。慢点在于状态复制和集合比较。位操作让包含、加入、删除、枚举子集变成常数或接近常数的机器指令。 如果面试中直接跳到优化而说不清暴力慢在哪里，后续复杂度分析就会显得像背模板。

\textbf{优化依据。} 异或解决成对抵消，按位计数处理出现三次，掩码表示访问集合或可用集合，低位提取用于枚举候选。 本题的优化要抓住“快慢指针不修改数组且常数空间；值域二分用抽屉原理。”这一点，把原来重复扫描或重复搜索的部分改成增量维护。优化以后，每次循环都应该只处理新增状态、删除过期状态或合并两个已经稳定的子答案，而不是重新审查全部输入。若题目条件发生变化，必须重新检查这个依据是否还成立。

\textbf{正确性不变量。} 掩码的每一位必须有固定含义。状态去重时，位置相同但掩码不同不能合并，因为未来能力不同。 用这道题复盘时，可以把不变量写成三句话：已经处理的部分满足什么，尚未处理的部分还保留什么可能性，当前操作为什么不会丢掉最优答案或合法答案。正确性来自二进制位独立性。异或按位满足交换结合和自反抵消；掩码 DP 则是对子集大小做归纳。 证明不需要很形式化，但必须能排除反例；如果你说“感觉这样选最好”，就还没有完成算法说明。

\textbf{C++ 实现要点。} 使用无符号整数能减少移位歧义。位数超过 31 时用 \code{long long} 或 \code{std::bitset}。表达式加括号，避免位运算优先级误读。本题还要特别注意：重复数出现多次、不能排序、不能改数组、长度为 n 加一。写代码时先把边界条件放在函数开头或循环不变量里，而不是用很多补丁式 if 修修补补。容器选择要和访问模式一致：需要顺序就别用无序表，需要极值就别完整排序，需要历史计数就别只保存布尔值。

\textbf{复杂度和边界。} 时间复杂度要按真实输入维度说明，不能把字符串长度、图边数、值域大小都混成一个 n。空间复杂度也要区分辅助结构和输出本身。测试时至少覆盖：最小输入、没有答案、答案在边界、重复元素或重复状态、以及会让暴力解最慢的规模输入。对于这题，边界复盘重点是：重复数出现多次、不能排序、不能改数组、长度为 n 加一。

\textbf{扩展和实验。} 快慢指针要求值都能作为合法下标。实验包含零、负数、最高位、重复数、全集掩码、空集合、只差一位的状态和资源耗尽后的不可达状态。实验报告建议记录三件事：暴力版本在什么规模开始变慢，优化版本减少了哪类重复工作，哪一个边界用例最容易打破错误实现。这样一张题卡才算从“看过题解”变成“能够迁移”。
